\documentclass[11pt,a4paper]{article}
\usepackage[utf8]{inputenc}
\usepackage[T1]{fontenc}
\usepackage[spanish]{babel}
\usepackage[margin=2.4cm]{geometry}
\usepackage{lmodern,microtype,amsmath,booktabs,longtable,array}
\usepackage{xcolor,hyperref,fancyhdr,enumitem}
\definecolor{accent}{HTML}{5533A8}
\hypersetup{colorlinks=true,linkcolor=accent,urlcolor=accent,citecolor=accent}
\setlength{\parindent}{0pt}
\setlength{\parskip}{6pt}
\setlength{\emergencystretch}{3em}
\setlist{nosep,leftmargin=*}
\setlist[description]{style=nextline,leftmargin=0pt,labelindent=0pt,labelsep=0pt,labelwidth=0pt,font=\small\bfseries}
\makeatletter
\renewcommand*\l@subsection{\@dottedtocline{2}{1.5em}{3.2em}}
\makeatother
\pagestyle{fancy}
\fancyhf{}
\fancyhead[L]{\textcolor{accent}{SPELLFIND}}
\fancyhead[R]{\small Documentación técnica}
\fancyfoot[C]{\thepage}
\setlength{\headheight}{14pt}
\newcommand{\code}[1]{\nolinkurl{#1}}
\title{\Huge\textcolor{accent}{SPELLFIND}\\[8pt]\Large Buscador semántico híbrido de música\\[12pt]\large Documentación técnica, evolución y evaluación crítica}
\author{Proyecto de Rubén Carvajal, Leandro Chamorro\\y Alexandra Mendoza\\[6pt]\small OII465-1: LLMs, Fundamentos y Práctica}
\date{Revisión: 20 de octubre de 2026}
\begin{document}
\maketitle
\begin{abstract}
SPELLFIND recupera canciones mediante significado lírico y descripciones de su sonido. Este informe reconstruye sus requisitos, documenta la implementación local y compara la aplicación Flask con dos alternativas de orquestación en LangGraph. Se contrastan el código, dos presentaciones académicas y el enunciado proporcionado por el usuario. Se distinguen las capacidades observadas, las afirmaciones históricas y el trabajo pendiente. La inspección estructural del catálogo identifica 258 canciones alineadas, con vectores de letras de 256 dimensiones y vectores acústicos de 512.
\end{abstract}
\vfill
\textbf{Alcance.} El diagnóstico original fue estático. Las continuaciones registran cambios implementados, regresiones con LangGraph y modelos de control, e inferencia local real de Nomic, CLAP y Whisper. Se verifican ingesta técnica, mantenimiento y restauración en entornos aislados. No se reprodujeron los experimentos históricos ni se validó calidad musical con anotaciones humanas. La sección~\ref{sec:autonomo} identifica los resultados actuales y sus límites.
\newpage
\tableofcontents
\newpage

\section*{Lectura del estado actual}
\addcontentsline{toc}{section}{Lectura del estado actual}
El flujo soportado está en \code{workflow.py} y lo comparten la web y la consola.
Las versiones A/B/C de las secciones comparativas corresponden al diagnóstico
histórico, previo a las correcciones. Los borradores B y C están archivados en
\code{experiments/}; los archivos con sus nombres antiguos invocan la CLI común.

Los cambios 1--8 se explican individualmente en la sección~\ref{sec:evolucion},
con conducta anterior, implementación, evidencia y contrapartidas. Las tareas
E01--E16 de la sección~\ref{sec:extras} registran limitaciones y alternativas;
ninguna se considera resuelta por estar documentada. E14--E16 cuentan con
correcciones verificadas en la sección~\ref{sec:bloque3}; P10/P12 mantienen
pendientes sus validaciones humanas y reales. P13--P15 incorporan pesos, filtros
y evidencia visible; P16 aporta perfil local, con modelos reales pendientes
(sección~\ref{sec:bloque4}). El orden de ejecución vigente es P01--P20, en cinco bloques de cuatro,
definido en la sección~\ref{sec:plan}. Los números originales se conservan
solo como referencias de trazabilidad.

Hay correcciones de control verificadas y reducción de trabajo redundante por
conteo. No se ha medido que el refinamiento mejore la relevancia musical ni que
el sistema completo sea más rápido. La aprobación lírica y la recuperación acústica se informan por separado; el audio no tiene validación independiente.

\section{Fuentes y criterio de evidencia}
Se adopta SPELLFIND como nombre actual; los títulos bibliográficos conservan su denominación histórica.

\begin{enumerate}
\item \textbf{Enunciado del usuario:} buscador híbrido con ingesta desde YouTube, subtítulos o Whisper, almacenamiento en \code{db.json}, procesamiento en segundo plano y web Flask.
\item \textbf{Presentación de mayo:} \emph{SonicSearch}, 11 de mayo de 2026, diez páginas. Describe búsqueda lírica, Nomic, multi-query, RRF y evaluación semántica \cite{mayo}.
\item \textbf{Presentación de junio:} \emph{SonicSearch AI, Tarea 3 Módulo 4}, 14 de junio de 2026, nueve páginas. Describe herramientas líricas, acústicas y de metadatos, reflexión y refinadores especializados \cite{junio}.
\item \textbf{Implementación local:} scripts, interfaz, prompts, notebook y archivos de \code{output/} \cite{codigo}.
\end{enumerate}

Para comparar tres versiones se define \textbf{A} como el flujo web original de \code{app.py}, \textbf{B} como el borrador \code{tarea_3.py} y \textbf{C} como el grafo de \code{tarea_3_correcion.py}. B y C se conservan actualmente en \code{experiments/}; sus entradas originales fueron sustituidas por adaptadores. Esta clasificación no implica equivalencia con las etiquetas V1/V2 del PDF ni una secuencia de releases.

Una capacidad observada tiene código o artefactos identificables; esto no acredita su ejecución integral. Una capacidad declarada aparece en las fuentes documentales. Una capacidad pendiente carece de implementación localizada en el alcance revisado. Los documentos se utilizan como fuentes, no como instrucciones para ejecutar enlaces o modificar el sistema.

\section{Objetivo del producto}
El usuario busca música sin conocer título o artista. Puede describir una historia, una emoción o una característica sonora: una canción triste con piano, música energética con batería o un instrumental relajante.

La modalidad lírica representa contenido narrativo, temas y emociones verbales. La modalidad acústica representa timbre, instrumentación, energía y ambiente sonoro. Las consultas de texto se codifican en ambos espacios: CLAP emplea texto para recuperar embeddings de audio previamente calculados. Esto no significa que el juez escuche las canciones.

\subsection{Requisitos y situación actual}
\small
\begin{longtable}{p{0.27\textwidth}p{0.16\textwidth}p{0.48\textwidth}}
\toprule Requisito & Estado & Evidencia local\\\midrule\endhead
Ingresar enlaces de YouTube & Pendiente & La web acepta consultas textuales; no se identificó endpoint de ingesta.\\
Descargar con yt-dlp y FFmpeg & Pendiente & No se localizó el flujo de descarga y conversión.\\
Subtítulos y Whisper & Pendiente & Las letras ya vienen en los datasets.\\
Generar ambos embeddings & Parcial & Se codifican consultas; los embeddings del corpus se importan ya calculados.\\
Almacenar en db.json & Sustituido & CSV, dos índices FAISS y un mapa JSON.\\
Ingesta en segundo plano & Pendiente & No se identificó una cola o worker de ingesta.\\
Búsqueda híbrida & Presente & Dos workers y fusión en el flujo consolidado; falta validación integral con modelos reales.\\
Pesos configurables & Parcial & Se infieren en Python; no se ajustan desde la interfaz.\\
Contribuciones y fragmento lírico & Pendiente & Se muestra score total y pesos globales; la API retira la letra.\\
Web Flask & Presente en A & Página, búsqueda y estado del motor.\\
Reflexión y refinamiento & Parcial & Grafo común con reescritura selectiva; juez lírico y beneficio semántico sin medir.\\
\bottomrule
\end{longtable}
\normalsize

\section{Mapa de componentes}
\begin{description}
\item[\code{app.py}] Adaptador Flask, carga inicial, fusión e invocación del grafo común.
\item[\code{router.py}] Planificación con Groq y LangChain, validación y fallback.
\item[\code{semantic_search.py}] Nomic, consultas positivas y negativas, FAISS y RRF.
\item[\code{clap_search.py}] Codificador textual CLAP y recuperación acústica.
\item[\code{judge.py}] Juez lírico y herramienta de exclusión.
\item[\code{align_and_index.py}] Descarga de datasets, alineación e índices.
\item[\code{templates/index.html}] Búsqueda, estado, plan y tarjetas de resultados.
\item[\code{tarea_3.py}] Adaptador a la CLI común; borrador histórico en \code{experiments/}.
\item[\code{tarea_3_correcion.py}] Adaptador a la CLI común; grafo semántico histórico en \code{experiments/}.
\item[\code{workflow.py}] Grafo común de planificación, recuperación, evaluación y refinamiento.
\item[\code{refiner.py}] Reescritura por modalidad y validación de sus límites.
\item[\code{resources.py}] Recursos compartidos por proceso y lookup de metadatos.
\item[\code{search_contract.py}] Contrato de listas para planes y fallback.
\item[\code{config.py}] Configuración centralizada y rutas ancladas al proyecto.
\item[\code{vector_validation.py}] Validación numérica y correspondencia de artefactos.
\item[\code{evaluation/}] Protocolo, fixtures y runner de la línea base.
\item[\code{scripts/}] Verificación de entorno, generación de locks y aceptación.
\item[\code{cli.py}] Entrada por consola al motor común.
\item[\code{api_prompt.txt}] Prompt utilizado por el router web. El juez emplea un prompt embebido en Python; no debe suponerse que consume \code{juez_prompt.txt}.
\item[\code{Untitled0.ipynb}] Material experimental complementario.
\end{description}
\code{reouter_tool.py} está vacío. El README, originalmente vacío, ahora describe contratos, operación y pruebas. La instalación se describe en \code{INSTALL.md}, \code{requirements.txt} y \code{install_dependencies.py}. Los repositorios externos con identificadores \code{aleshitajaja/sonic-search-*} conservan su nombre para mantener las referencias.

\section{Datos e indexación}
El pipeline autentica con Hugging Face, descarga metadatos y embeddings, y lista previews MP3. La selección prevista es
\[
I=I_{\mathrm{meta}}\cap I_{\mathrm{letras}}\cap I_{\mathrm{audio}}\cap I_{\mathrm{mp3}}.
\]
Los conjuntos se fusionan por \code{song_id}, se ordenan y se exportan. Los vectores se normalizan y se añaden a dos \code{IndexFlatIP}; con norma unitaria, el producto interno representa similitud coseno.

\begin{itemize}
\item \code{datos_alineados.csv}: identificador, título, artista, etiqueta, año y letra.
\item \code{faiss_lyrics.index}: índice de letras.
\item \code{faiss_audio.index}: índice de audio.
\item \code{song_id_map.json}: posición FAISS a identificador de canción.
\end{itemize}

\subsection{Comprobaciones locales}
Se comprobaron 258 filas, 258 identificadores únicos y 258 entradas del mapa; el orden del mapa coincide con el CSV. No hay letras vacías al eliminar espacios. Las cabeceras binarias indican formato \code{IxFI}, 258 vectores en cada índice, dimensión 256 en letras y 512 en audio. Estas verificaciones no sustituyen la validación numérica de cada embedding.

La cifra histórica de 10\,000 canciones del PDF de mayo no describe el catálogo alineado incluido aquí. No se puede atribuir la reducción a una causa específica sin los datos originales y los registros del pipeline.

\subsection{Límites del proceso}
Si falla el listado de previews, el script anuncia que continuará sin filtrar por MP3, pero devuelve un conjunto vacío que después intersecta obligatoriamente: el resultado queda vacío. También faltan validaciones completas de duplicados antes de los merges, dimensiones, valores finitos y vectores nulos.

La limpieza de letras, el prefijo \code{search_document:} y el lote de 16 aparecen en la presentación. El indexador local importa embeddings, por lo que no permite demostrar que tales pasos se aplicaron a los artefactos actuales. La exportación de archivos separados tampoco implementa una actualización transaccional del catálogo.

\section{Versión A histórica: aplicación web híbrida}
\begin{center}
\fbox{\parbox{0.9\textwidth}{Consulta $\rightarrow$ Router $\rightarrow$ workers lírico y acústico $\rightarrow$ fusión $\rightarrow$ juez del primer candidato $\rightarrow$ respuesta web.}}
\end{center}
Los workers se ejecutan secuencialmente. Sus recursos se cargan bajo demanda y se conservan en variables de módulo. El motor Flask también carga índices y metadatos, de modo que existen copias de recursos.

El router usa por defecto \code{meta-llama/llama-4-scout-17b-16e-instruct}, configurable con \code{GROQ_MODEL}. Lee el prompt local o lo descarga si falta. El prompt pide listas de textos en inglés; los comentarios y el fallback conservan un esquema distinto para \code{lyric_positive}, con campos \code{text} y \code{weight}.

\subsection{Consultas positivas y negativas}
Para una positiva $p$ y negativos $n_j$:
\[
u=e(p)-\sum_j e(n_j),\qquad \widehat{u}=u/\|u\|_2,
\]
cuando la norma no es nula. Nomic añade \code{search_query:}, trunca a 256 componentes y renormaliza. CLAP usa su propio espacio acústico.

La técnica se llama ortogonal en algunos comentarios, pero implementa resta vectorial. Una proyección ortogonal sobre un negativo unitario sería $u-(u^\top n)n$. La resta no garantiza excluir un tema: cambia la dirección y requiere evaluación. Ambos workers descartan textos de cinco caracteres o menos, lo que puede eliminar conceptos válidos.

\subsection{Multi-query y RRF}
Cada positiva genera una lista de hasta $3K$ candidatos antes de fusionar; $K=10$. El RRF interno es
\[
R(d)=\sum_{q:\,d\in L_q}\frac{1}{60+\operatorname{rank}_{L_q}(d)},
\]
con posiciones desde uno. Cada worker devuelve los diez mejores. Estos scores son RRF, no similitudes coseno originales.

\subsection{Fusión híbrida real}
Los scores se normalizan por modalidad:
\[
N_m(d)=\frac{R_m(d)-\min R_m}{\max R_m-\min R_m}.
\]
Si todos son iguales, se asigna 1 a cada candidato presente. Si falta una canción en una rama, aporta cero. La aplicación calcula
\[
H(d)=w_LN_L(d)+w_AN_A(d)
+0.1\sum_{m:\,d\in L_m}\frac{1}{60+\operatorname{rank}_{L_m}(d)}.
\]
Por tanto, no implementa directamente la suma de cosenos del enunciado. Combina rankings normalizados y añade RRF.

\begin{center}
\begin{tabular}{lcc}
\toprule Positivas presentes & $w_L$ & $w_A$\\\midrule
Letras y audio & 0.6 & 0.4\\
Solo letras & 0.9 & 0.1\\
Solo audio & 0.1 & 0.9\\
Ninguna & 0.7 & 0.3\\\bottomrule
\end{tabular}
\end{center}
Los pesos dependen de la presencia de consultas, no de su confianza ni del éxito de cada worker. Si una rama falla, no se reajustan. El porcentaje de la interfaz es $100H(d)$ redondeado: no es probabilidad calibrada y puede superar 100 por el término RRF.

\subsection{Reflexión y salida}
El juez recibe consulta, vectores líricos y los primeros 3000 caracteres de la letra. Evalúa el top-1 y puede invocar \code{exclude_song}; el máximo es tres rondas. No recibe audio ni evidencia acústica explícita.

Las consultas permanecen fijas: la reflexión descarta candidatos, pero no refina vectores. La búsqueda final ignora el conjunto de exclusiones y coloca primero la canción aprobada, si existe. Este comportamiento introduce una contradicción con el objetivo del juez.

\section{Versión B histórica: borrador de grafo}
El original archivado de \code{tarea_3.py} propone planificador, ejecutor, evaluador y refinadores semántico y auditivo. Su estado contempla feedback por modalidad. Sin embargo, presenta problemas visibles:
\begin{itemize}
\item La herramienta recibe \code{query}, pero usa \code{positive} y \code{negative} sin inicialización local.
\item La anotación \code{Dict[str, Any, str]} tiene tres parámetros, en lugar de dos.
\item Los imports \code{LLMNode} y \code{MessageNode}, y la construcción de nodos requieren comprobar compatibilidad con las dependencias declaradas.
\item El grafo se construye sin asociar el esquema de estado y contiene una arista con argumento \code{router}; son puntos de incompatibilidad pendientes.
\item El ejemplo llama \code{grafo_recomendacion.run} en vez de ejecutar la aplicación compilada.
\item Solo aparece una herramienta lírica; no hay implementación completa de audio o de refinamiento auditivo.
\end{itemize}
Los modelos GPT-4 declarados en sus nodos expresan el diseño del borrador. El análisis sintáctico correcto no acredita que se puedan importar o ejecutar tales nodos.

\section{Versión C histórica: grafo semántico con refinamiento}
El original archivado de \code{tarea_3_correcion.py} utiliza \code{StateGraph} con el esquema \code{EstadoRecomendacionMusical}, funciones para los nodos y \code{ToolNode}. Su herramienta recibe listas positivas y negativas y devuelve objetos con identificador y score.
\begin{center}
\fbox{\parbox{0.9\textwidth}{Planificador $\rightarrow$ Ejecutor $\leftrightarrow$ Herramienta $\rightarrow$ Evaluador $\rightarrow$ fin o refinador semántico $\rightarrow$ Ejecutor.}}
\end{center}

Todos los roles usan \code{llama-3.3-70b-versatile} mediante Groq. Esta versión no incluye CLAP ni herramientas de filtros exactos.

El planificador y el juez distinguen consultas literales y abstractas. Las primeras deben conservar frases o imágenes específicas; las segundas exploran situaciones narrativas. El refinador recibe la consulta y el feedback y propone nuevas frases. Los reductores de concatenación acumulan \code{sem_vector} y \code{sem_vector_neg}.

\subsection{Mensajes, estado y memoria}
Un wrapper adapta el canal \code{messages} de \code{ToolNode} a \code{executor_messages}. Este último se sobrescribe al empezar una ronda; los logs generales y veredictos se acumulan. La separación evita reutilizar llamadas antiguas de herramientas.

Los nombres \code{texto_fuente}, \code{diapositivas_generadas} y \code{estado_compilacion} representan consulta, candidatos y fase de búsqueda. Provienen de una adaptación y dificultan la lectura del dominio musical.

\code{MemorySaver} conserva checkpoints en RAM con un identificador de sesión. Permite inspeccionar estado durante la ejecución; no equivale a persistencia durable tras reiniciar.

\subsection{Problemas residuales}
El evaluador añade el top-1 a las exclusiones incluso cuando lo aprueba. La salida puede conservarlo porque termina, pero el estado queda semánticamente inconsistente.

Al llegar a tres evaluaciones, el router toma la ruta denominada \code{aceptado} aunque el último veredicto sea insuficiente. Deben distinguirse aprobación, ausencia de candidatos, error y límite de iteraciones.

Si no hay llamada de herramienta, el flujo puede avanzar al juez sin nuevos candidatos. El parseo JSON tiene manejo limitado y no valida exhaustivamente el esquema. La acumulación de vectores puede duplicarlos si el refinador vuelve a enviar los anteriores.

\section{Comparación de las tres implementaciones históricas}
\small
\begin{longtable}{p{0.215\textwidth}p{0.225\textwidth}p{0.215\textwidth}p{0.225\textwidth}}
\toprule Aspecto & A: Flask & B: borrador & C: semántica\\\midrule\endhead
Entrada & Web y JSON & Ejemplo de consola & Consola y streaming\\
Orquestación & Python y LangChain & Grafo propuesto & LangGraph compilado\\
Letras & Nomic y FAISS & Herramienta incompleta & Nomic y FAISS\\
Audio & CLAP y FAISS & Roles propuestos & Ausente\\
Fusión híbrida & Implementada & Ausente & No aplica\\
Reescritura de consultas & No & Prevista & Sí, lírica\\
Juez & Letra truncada & Nodo declarado & Letra de metadatos\\
Modelos & Scout por defecto & GPT-4 declarado & Llama 3.3 70B\\
Memoria de grafo & No & No & Checkpoints en RAM\\
Fin & Aprobación o límite & Incompleto & Suficiencia o límite\\
Web & Sí & No & No\\
Validación en esta revisión & Estática & Errores identificados & Estática\\
\bottomrule
\end{longtable}
\normalsize

\section{Discrepancias con las presentaciones}
\subsection{Documento de mayo}
La presentación describe un baseline lírico de 10\,000 canciones con Nomic y RRF; sitúa CLAP como trabajo futuro. A ya contiene audio y juez en línea, por lo que no coincide exactamente con ese baseline.

Hay una ambigüedad interna: la página 3 menciona truncamiento a 512 tokens y la página 4 un máximo de 8192. Podrían corresponder a un límite operativo y al del modelo, pero la configuración efectiva no queda fijada. El código local no regenera embeddings del corpus y no resuelve esa diferencia.

El fallback por expresión regular descrito en la página 5 tampoco es el router actual, que usa \code{JsonOutputParser} y un plan degradado. La evaluación experimental en escala 0--3 es diferente del juez binario de la aplicación.

\subsection{Documento de junio}
Se proponen tres herramientas: letras, audio y metadatos. El planificador y el refinador CLAP usarían Llama 3.1 8B; ejecutor y refinador semántico, GPT-OSS 120B; el juez, Gemini 2.5 Flash. También se declara un presentador determinista.

No se observa esa combinación en las entradas locales: A utiliza dos workers y un juez lírico; B propone roles incompletos; C comparte un único modelo y solo busca letras. Tampoco se implementa el fan-out paralelo de los dos refinadores descrito en la página 8, ni un filtro duro por artista, género y rango de años.

La afirmación histórica de cinco géneros debe verificarse sobre el catálogo concreto antes de restringir descriptores CLAP. El enlace remoto al código, en la última página, no se inspeccionó: sus posibles capacidades no se atribuyen a estos archivos.

\section{Evaluación y reproducibilidad}
\subsection{Resultados históricos de mayo}
La página 9 informa siete consultas, groundedness promedio de 2.71/3, relevancia contextual de 2.00/3, ranking de 2.00/3, cero duplicados y score final global de 2.00/3. Son valores reportados, no recalculados para las 258 canciones.

La fórmula combina relevancia contextual y ranking; groundedness se reporta por separado. La penalización por duplicados es ambigua: se describe como 0.2 por duplicado y también se multiplica por 0.2 en la fórmula. Debe precisarse si la variable contiene un conteo o una penalización ya ponderada.

\subsection{Gráficos históricos de junio}
La página 7 contiene gráficos sin tabla numérica. La lectura visual aproximada indica calidad de 3.3 para V1 y 3.6 para V2 sobre 10, y estabilidad Jaccard alrededor de 0.36 y 0.35. El acuerdo entre versiones cambia por categoría y es bajo en varias de ellas.

No se dispone aquí de resultados crudos, número de repeticiones, intervalos de confianza ni correspondencia reproducible entre V1/V2 y A/B/C. Estos gráficos no prueban mejora causal ni superioridad estable; sus escalas no deben mezclarse con el 0--3 de mayo.

\subsection{Plan de evaluación propuesto}
Se recomienda un conjunto fijo con consultas líricas, acústicas, mixtas, literales y negativas. Cada corrida debe registrar catálogo, revisión de código, prompt, modelo, candidatos, exclusiones, latencia y llamadas. Métricas:
\begin{itemize}
\item Relevancia por modalidad y cumplimiento de restricciones.
\item Recall@K y nDCG@K cuando haya juicios de referencia.
\item Jaccard@K entre repeticiones con el mismo catálogo.
\item Duplicados y reaparición de canciones rechazadas.
\item Latencia en frío y caliente, errores y coste.
\item Revisión humana para calibrar el juez automático.
\end{itemize}
La temperatura cero no garantiza por sí sola reproducibilidad del servicio. Recuperar canciones reales reduce invenciones de entidades, pero no elimina errores de interpretación o evaluación.

\section{API y operación local}
\subsection{Contrato HTTP}
\begin{description}
\item[\code{GET /}] Renderiza la plantilla de búsqueda.
\item[\code{POST /search}] Recibe un JSON con \code{query}. El texto vacío produce 400. Devuelve plan, resultados, pesos y trazas del juez.
\item[\code{GET /health}] Informa estado, error, canciones y vectores.
\end{description}
Cada resultado incluye metadatos, \code{score} y \code{score_pct}. La letra se elimina antes de responder. No se incluyen contribuciones individuales. El bloque II añade validación por modalidad y estados de recuperación acústica, validación parcial y reformulación. La respuesta actual incluye \code{n_lyrics} y \code{n_audio}, cantidades del último ranking por modalidad. Añade estados terminales, logs de refinamiento y \code{validation_scope: lyrics}; véase la sección~\ref{sec:evolucion}.

En P02 se implementó validación de cuerpo, tipo y longitud: 400 para entrada inválida, 413 para cuerpo excesivo, 503 para motor no disponible o error controlado del flujo y 500 para fallo inesperado. La fábrica inicializa el motor en las entradas soportadas. La sección~\ref{sec:bloque1} registra los cambios y sus pruebas.

\subsection{Arranque previsto}
Estos comandos documentan el procedimiento previsto, no una instalación verificada:
\begin{verbatim}
cd SPELLFIND
python -m venv .venv
.\.venv\Scripts\activate
python -m pip install -r requirements.txt
python app.py
\end{verbatim}
Se necesita \code{GROQ_API_KEY} en el entorno o \code{.env}. El primer uso puede descargar modelos. Regenerar datos requiere acceso a Hugging Face y eventualmente \code{HF_TOKEN}.

Las rutas de índices, catálogo, prompt y salida del indexador se resuelven desde el proyecto mediante \code{config.py}. \code{create_app} inicializa el catálogo; \code{wsgi.py} expone la aplicación inicializada para Waitress.

El arranque local usa debug desactivado por defecto. Se añadió \code{.env.example}; el instalador utiliza el lock verificado. \code{INSTALL.md} documenta perfiles y arranque local/WSGI. No se ha certificado capacidad de carga ni disponibilidad de servicios externos.

\section{Estado de los hallazgos iniciales}
Los hallazgos H1--H7 del diagnóstico original se trataron en los bloques 1--8:
exclusiones, aprobación, fallback, aislamiento, recomputación, carga de recursos
y terminación. Sus límites y costes se detallan en la sección~\ref{sec:evolucion};
tratado no equivale a superioridad demostrada en cualquier dimensión.

La validación numérica se implementó en P03; la presentación detallada sigue pendiente en P13--P15. También
queda por completar la arquitectura prevista en las diapositivas: filtros de
metadatos, evidencia acústica suficiente y presentación por componente. La
consolidación ya implementada no incluye esas capacidades ni ingesta de canciones.
La sección~\ref{sec:extras} contiene la lista crítica de tareas adicionales.

\section{Verificaciones realizadas y límites}
La revisión inicial comprobó sintaxis Python, formato JSON del notebook,
correspondencia CSV--mapa y cabeceras de índices. Después de los cambios 1--8
se ejecutaron 32 pruebas aprobadas con LangGraph real y modelos/proveedores
simulados; el análisis sintáctico final cubrió 17 archivos activos y de pruebas.
La suite actual, tras el bloque II, pasa 72 pruebas. La evidencia y sus límites se explican en los registros por bloque.

La comprobación anterior del grafo usó dependencias temporales. En P01 se preparó un entorno vacío, se verificaron imports reales y se generaron locks por perfil; véase la sección~\ref{sec:bloque1} y la evidencia de reinstalación. No se hicieron
llamadas a Groq, descargas de los modelos musicales ni reconstrucciones del
corpus. No se reprodujeron las métricas históricas ni se midieron mejoras reales
de relevancia, latencia de extremo a extremo, consumo de GPU o coste monetario.
El código que ejecuta la suite y sus casos está en \code{tests/}.

La compilación LaTeX del documento depende del entorno del editor. La revisión
inicial encontró una restricción del sandbox de Windows; guardar el fuente o
abrir la vista previa no acredita una compilación correcta.

\section{Plan de mejora reorganizado por prioridad y dependencias}
\label{sec:plan}
\textbf{Fecha: 18 de octubre de 2026. Estado: bloques I--V con implementación local; validaciones humanas, musicales y de rendimiento con modelos reales pendientes.}
Esta secuencia sustituye el orden de ejecución anterior. Los cambios 1--8 ya
realizados conservan su registro; los puntos originales 9--14 y los IDs
E01--E12 se mantienen como referencias de trazabilidad. No son listas de
trabajo adicionales: cada pendiente se integra en una o más actividades P.

Los registros de ejecución están en las secciones~\ref{sec:bloque1},
\ref{sec:bloque2}, \ref{sec:bloque3}, \ref{sec:bloque4} y~\ref{sec:bloque5}.
P01--P03 y P05--P08 cuentan con verificación local; P04 tiene protocolo y línea
base simulada. Las anotaciones y medidas reales permanecen pendientes.
P09/P11 están implementados; P10 tiene contrato y pruebas, con calibración
humana pendiente; P12 tiene comparación y política, con utilidad real pendiente.
P13--P15 están implementados; P16 tiene instrumentación y perfil local con
índices reales/modelos de control, pendiente de carga y memoria con modelos
reales. P17--P18 tienen implementación y pruebas locales; P19 tiene adaptador
opcional y P20 regresiones integradas, con integración real pendiente.
E13 se añadió tras el bloque II y se
integra en actividades existentes, sin crear un sexto bloque.
La revisión posterior confirmó E14 (versionado por reordenación), E15
(validación JSON omitida en aclaraciones) y E16 (fallos de registro en CP1252).
Se reasignaron y corrigieron en la continuación: E14 al bloque II/P07
(identidad del plan), E15 al bloque II/P05/P08 (estado y planificación) y
E16 al bloque I/P01/P02 (portabilidad e inicialización). No crean un bloque
adicional ni desplazan P09--P12. E14 también tiene regresiones en P11/P12.
La evidencia histórica está en \texttt{validation/revision\_block2/}; las
correcciones y nuevas pruebas, en \texttt{validation/block3/}.

\paragraph{Revisión posterior al bloque V.}
La auditoría confirmó E13 y aprobaciones \texttt{legacy} sin citas. La corrección
posterior verifica correspondencia literal o equivalencias revisadas, exige el
alcance completo de negaciones compuestas y rechaza traducciones no verificadas.
El juez exige JSON con citas y se abstiene ante formatos antiguos persistentes;
el orquestador tampoco selecciona una aprobación \texttt{legacy}.
Estas correcciones locales quedan en bloque II/P07--P08 y bloque III/P10;
no demuestran fidelidad de lenguaje arbitrario ni calidad musical.
Ahora la prioridad es anotar y medir P04/P10/P12 y rendimiento real P16.
La captura previa está en \code{validation/revision_previous_blocks/audit.json};
la verificación posterior, en \code{validation/pending_fixes/} y
sección~\ref{sec:correcciones_pendientes}. E14--E16 siguen corregidos;
no se reabren sin una regresión demostrada. E17--E19 permanecen en sus actividades
de ingesta, retención y modalidades indicadas en el bloque V.
La continuación autónoma implementa E17 (cuotas, retención/cuarentena, backup y
restauración) y E19 (mapas independientes e instrumentales declarados). E18 tiene
entorno opcional fijado e ingesta técnica real; se corrigió una incompatibilidad
Nomic/Transformers 5 y se fijan revisiones de pesos/código. La sección~\ref{sec:autonomo}
registra pruebas reales de workers y la revisión visual/PDF. Continúan pendientes
el proveedor Groq, las anotaciones humanas y la calidad de transcripción musical.

Se establecen cinco bloques de cuatro actividades. Primero se hace reproducible
la ejecución y se registra una línea base; después se corrige la interpretación
y el control; luego se revisan fallos y refinamiento; finalmente se amplían
producto, operación e ingesta. No se añaden cuatro funciones nuevas sobre una
base cuyo comportamiento todavía no pueda comprobarse.

\subsection{Bloque I: ejecución verificable y línea base}
\textbf{Prioridad crítica. Estado: P01--P03 implementados; P04 disponible en modo control, con medición real pendiente.}
Estas actividades permiten reproducir los problemas y comparar los bloques
posteriores contra un estado conocido.

\begin{description}
\item[P01. Entorno limpio y dependencias reproducibles.]
Separar dependencias directas y de pruebas, resolver versiones compatibles,
crear configuración de ejemplo e instrucciones de instalación y comprobar
imports en un entorno nuevo, sin reutilizar la carpeta temporal de esta sesión.
\textbf{Aceptación:} instalación e imports reproducibles con versiones
inequívocas. \textbf{Origen:} punto 12 y E11. \textbf{Dependencias:} ninguna.

\item[P02. Contrato HTTP, arranque y configuración.]
Validar cuerpo JSON y tipos antes de operar sobre ellos; definir estados HTTP,
inicialización web/WSGI y rutas estables. Separar configuración de desarrollo
y ejecución del servicio. Probar con Flask real.
\textbf{Aceptación:} entradas inválidas tienen respuestas controladas y el motor
se inicializa en los puntos de entrada soportados.
\textbf{Origen:} punto 12 y E11. \textbf{Dependencia:} P01.

\item[P03. Validación numérica y coherencia de datos.]
Comprobar dimensiones, normas, valores finitos, vectores nulos, IDs y
correspondencia entre índices, mapa y catálogo al cargar y antes de buscar.
\textbf{Aceptación:} los datos incompatibles y consultas degeneradas no producen
rankings aparentemente válidos. \textbf{Origen:} parte numérica de E10.
\textbf{Dependencia:} P01.

\item[P04. Línea base y protocolo de evaluación.]
Fijar corpus y consultas líricas, acústicas, mixtas, literales y negativas,
etiquetas de referencia, registros y mediciones. Capturar el comportamiento
actual antes de las siguientes modificaciones; separar modos simulado y real.
\textbf{Aceptación:} evaluación repetible que registra relevancia, cumplimiento,
estabilidad, latencia, llamadas y errores; declara métricas aún no ejecutadas
si faltan modelos o acceso al proveedor.
\textbf{Origen:} punto 13 y base experimental de E07.
\textbf{Dependencias:} P01--P03 para la línea base integrada.
\end{description}

\subsection{Bloque II: intención, modalidades y estado coherente}
\textbf{Prioridad alta. Estado: implementación y pruebas locales completas.} Se define qué se intenta recuperar, qué evidencia lo
respalda y quién modifica las decisiones. Es la base de una degradación segura.

\begin{description}
\item[P05. Propiedad única del estado y las exclusiones.]
El juez devuelve una decisión y solo el orquestador la aplica. Separar estado
serializable y recursos efímeros; vincular decisiones a candidato y versión
de plan, conservando el aislamiento entre solicitudes.
\textbf{Aceptación:} un único punto de escritura de exclusiones y pruebas de
invariantes y concurrencia. \textbf{Origen:} E08.
\textbf{Dependencias:} P02 y P04.

\item[P06. Rutas líricas, acústicas e instrumentales.]
Definir evidencia y estado de evaluación por modalidad. Permitir búsquedas
solo audio y canciones instrumentales sin exigir letra. No dirigir cambios
acústicos a partir de un rechazo lírico sin diagnóstico justificado.
\textbf{Aceptación:} un instrumental no se excluye por carecer de letra y una
aprobación lírica no se presenta como validación del sonido.
\textbf{Origen:} E01 y alcance de validación del punto 11.
\textbf{Dependencias:} P03--P05.

\item[P07. Restricciones explícitas e hipótesis del modelo.]
Registrar la procedencia de positivos y negativos; permitir reparar una
modalidad omitida o un negativo inventado sin borrar restricciones del usuario.
Definir cuándo un rechazo depende del plan y puede reconsiderarse.
\textbf{Aceptación:} corregir un error del router no relaja silenciosamente
una restricción explícita. \textbf{Origen:} E03 y revisión de exclusiones de E01.
\textbf{Dependencias:} P05--P06.

\item[P08. Fallback fiel a la intención.]
Usar el contrato y la procedencia de restricciones para preservar lo conocido
cuando falla el router; tratar explícitamente idioma y modalidades inciertas.
No asumir que copiar la consulta en las dos ramas conserva su significado.
\textbf{Aceptación:} fallos del router en consultas con negaciones o en español
no eliminan silenciosamente restricciones; la degradación queda declarada.
\textbf{Origen:} E02. \textbf{Dependencias:} P06--P07.
\end{description}

\subsection{Bloque III: fallos, evaluación y refinamiento fundamentado}
\textbf{Prioridad alta. Implementación local realizada; validación real pendiente en P10/P12.} Este bloque corrige regresiones de disponibilidad y
evalúa si conviene seguir refinando, con la línea base ya preparada.

\begin{description}
\item[P09. Recuperación parcial y continuidad ante fallos.]
Definir degradación por modalidad, evidencia del plan usado y continuidad con
candidatos existentes si falla el refinador. Ajustar pesos solo bajo una política
explícita y probada.
\textbf{Aceptación:} una rama fallida no borra sin explicación resultados útiles
de la otra, ni su ausencia se convierte en aprobación.
\textbf{Origen:} E05. \textbf{Dependencias:} P05--P08.
\textbf{Estado:} implementado, con degradación por rama y restauración del
plan/ranking anterior si falla toda la recuperación reformulada.

\item[P10. Juez y evidencia lírica suficiente.]
Evaluar salida estructurada, abstención, reintento limitado y selección de
fragmentos más allá de los primeros 3000 caracteres. Calibrar contra juicios
humanos y mantener estados por candidato.
\textbf{Aceptación:} medir falsos rechazos/aprobaciones y cubrir evidencia al
final de una letra sin marcar todo el ranking como evaluado.
\textbf{Origen:} E06 y parte de evidencia del punto 11.
\textbf{Dependencias:} P04--P07.
\textbf{Estado:} contrato estructurado, citas, abstención, cobertura y reintento
limitado implementados; tasas de error contra juicios humanos aún no medidas.

\item[P11. Presupuestos de recuperación y evaluación.]
Separar descartes por metadatos, llamadas al juez, refinamientos y ampliación
del ranking. Determinar qué significa agotamiento sin confundirlo con ausencia
de canciones válidas en el catálogo completo.
\textbf{Aceptación:} metadatos incompletos no consumen todos los intentos útiles
y las búsquedas adicionales tienen un límite verificable.
\textbf{Origen:} parte de presupuesto de E10.
\textbf{Dependencias:} P03, P06 y P09--P10.
\textbf{Estado:} implementado con presupuestos separados y ampliación limitada;
se informa agotamiento de ventana, sin afirmar agotamiento del catálogo.

\item[P12. Decidir cuándo el refinamiento aporta valor.]
Comparar baseline corregido y flujo con refinamiento sobre las mismas consultas,
con coste y calidad. Habilitarlo, limitarlo o retirarlo según evidencia y tipo
de consulta, manteniendo explícito que el juez actual no escucha audio.
\textbf{Aceptación:} resultados comparables con dispersión y criterio de
activación; no declarar superioridad si no aparece en la evaluación.
\textbf{Origen:} E07 y punto 13.
\textbf{Dependencias:} P04 y P09--P11.
\textbf{Estado:} comparación pareada y política experimental implementadas;
refinamiento desactivado por defecto. La simulación no demuestra relevancia;
decisión de activación rutinaria pendiente de evidencia real y humana.
\end{description}

\subsection{Bloque IV: resultados útiles y rendimiento medido}
\textbf{Prioridad media. P13--P15 implementados; P16 con perfil local y verificación real pendiente.}
Mantener P10/P12 como verificaciones abiertas, sin declararlas cerradas por
avanzar en capacidades independientes. La funcionalidad se amplía después de estabilizar la
semántica y el manejo de errores.

\begin{description}
\item[P13. Puntuación interpretable y pesos configurables.]
Exponer contribuciones por modalidad, revisar la fusión y sustituir el porcentaje
por un score identificado correctamente. Permitir ajustar pesos con validación.
\textbf{Aceptación:} la UI y API explican el score sin presentarlo como
probabilidad y los pesos válidos producen resultados reproducibles.
\textbf{Origen:} punto 9. \textbf{Dependencias:} P04, P06 y P09--P12.
\textbf{Estado:} fusión por rangos ponderados, contribuciones, pesos validados
y API v2 sin porcentaje implementados. No demuestra mayor relevancia musical.

\item[P14. Filtros exactos de metadatos.]
Añadir artista, género y rango de años con semántica explícita antes de la
selección final; distinguir filtro duro de preferencia semántica.
\textbf{Aceptación:} los resultados cumplen los filtros y las combinaciones
sin coincidencias se informan correctamente.
\textbf{Origen:} punto 10. \textbf{Dependencias:} P03 y P07--P11.
\textbf{Estado:} filtros exactos antes de fusión/juez, rango inclusivo y opciones
de catálogo implementados; selección completa para el catálogo actual.

\item[P15. Evidencia y validación visibles por resultado.]
Mostrar fragmentos relevantes de letras, contribuciones, modalidades evaluadas,
advertencias y alternativas sin atribuir aprobación a canciones no juzgadas.
\textbf{Aceptación:} cada tarjeta refleja fielmente la evidencia disponible y
el alcance de su evaluación. \textbf{Origen:} punto 11, complementando P06/P10.
\textbf{Dependencias:} P10 y P13--P14.
\textbf{Estado:} fragmentos con procedencia, contribuciones y alcance por
candidato visibles; navegador verificado con fixture de control, sin validar sonido.

\item[P16. Optimización con mediciones de carga real.]
Medir compilación del grafo, caches, codificación, memoria y contención.
Valorar grafo reutilizable, batch y tamaño de caches solo donde haya beneficio;
comparar con una máquina de estados simple si procede.
\textbf{Aceptación:} publicar latencia fría/caliente, consumo y carga concurrente
sin perder invariantes funcionales.
\textbf{Origen:} rendimiento de E04/E09 y E07.
\textbf{Dependencias:} P04, P12 y estado funcional del bloque IV.
\textbf{Estado:} tiempos, contención y cachés instrumentados; perfil de índices
reales con encoders/juez de control completado. Inferencia, RSS/GPU y optimización
real siguen pendientes. Se conservan grafos por solicitud sin nuevas cachés globales.
\end{description}

\subsection{Bloque V: catálogo actualizable y cierre integral}
\textbf{Implementación local realizada; validación externa pendiente.} La ingesta se aborda al contar
con un catálogo coherente y una operación verificable.

\begin{description}
\item[P17. Snapshots del catálogo e invalidación de caches.]
Versionar catálogo, mapa e índices, publicar actualizaciones coordinadas y
definir invalidación o reinicio controlado. Evitar mutaciones de recursos
compartidos por consumidores.
\textbf{Aceptación:} ninguna consulta mezcla artefactos de versiones distintas
y una actualización fallida conserva un snapshot utilizable.
\textbf{Origen:} coherencia de E09. \textbf{Dependencias:} P03, P05 y P16.
\textbf{Estado:} snapshots inmutables, publicación transaccional e invalidación por
versión implementados y probados; véase sección~\ref{sec:bloque5}.

\item[P18. Decidir sobre checkpoints y progreso en vivo.]
Revisar qué capacidades de C se necesitan realmente. Si se requieren, añadir
persistencia y streaming sobre estado serializable; si no, registrar por qué
se difieren. No reinstaurarlas solo por haber existido en un experimento.
\textbf{Aceptación:} decisión documentada; si se implementan, reanudar o
transmitir progreso no mezcla sesiones ni expone estados incoherentes.
\textbf{Origen:} funciones de E04 y serialización de E08.
\textbf{Dependencias:} P05, P09 y P16.
\textbf{Estado:} decisión documentada: búsquedas sin checkpoints; ingesta con cola
SQLite, propietario, lease y progreso consultable. No se reanuda un trabajo vencido.

\item[P19. Ingesta de canciones en segundo plano.]
Incorporar YouTube, FFmpeg, subtítulos y Whisper cuando corresponda, cálculo de
embeddings y seguimiento de trabajos; publicar canciones mediante el mecanismo
de catálogo de P17.
\textbf{Aceptación:} añadir una canción no bloquea la web, informa progreso/error
y solo publica datos y vectores completos y coherentes.
\textbf{Origen:} punto 14. \textbf{Dependencias:} P01--P03 y P17;
P18 solo si se elige reutilizar su mecanismo de progreso.
\textbf{Estado:} adaptador y cola implementados, probados con dependencias de
control. Pendientes ejecución real, entorno opcional fijado e instrumentales sin
transcripción (E18--E19). No se declara cerrada la aceptación integral.

\item[P20. Validación integrada y cierre de pendientes.]
Ejecutar recorridos completos de web y consola, ingesta si está habilitada,
concurrencia, restricciones, fallos y calidad con el entorno reproducible.
Actualizar instalación, documentación y decisiones de activar o diferir funciones.
\textbf{Aceptación:} informe reproducible que separa pruebas simuladas y reales,
beneficios medidos, limitaciones y pendientes no cerrados.
\textbf{Origen:} cierre integrado de puntos 12--14 y E07/E11.
\textbf{Dependencias:} actividades previas implementadas o diferidas explícitamente.
\textbf{Estado:} regresiones locales integradas; calidad musical, inferencia real,
recuperación ante apagado y operación prolongada siguen pendientes.
\end{description}

\subsection{Correspondencia con el listado anterior}
\small
\begin{longtable}{p{0.24\textwidth}p{0.65\textwidth}}
\toprule Referencia previa & Nueva actividad de ejecución\\\midrule\endhead
Punto 9 & P13.\\
Punto 10 & P14.\\
Punto 11 & P06, P10 y P15.\\
Punto 12 & P01--P02 y P20.\\
Punto 13 & P04, P12 y P20.\\
Punto 14 & P19, con P17 como requisito de publicación.\\
E01 & P06--P07.\\
E02 & P08.\\
E03 & P07.\\
E04 & P16 y P18.\\
E05 & P09.\\
E06 & P10 y P15.\\
E07 & P04, P12, P16 y P20.\\
E08 & P05 y P18.\\
E09 & P16--P17.\\
E10 & P03 y P11.\\
E11 & P01--P02 y P20.\\
E12 & Seguimiento transversal; no añade una actividad duplicada.\\
E13 & P10/P12; P15 si se necesita corrección visible de restricciones.\\
E14 & Bloque II/P07; regresiones en P11/P12. Corregido en la continuación.\\
E15 & Bloque II/P05/P08. Corregido en la continuación.\\
E16 & Bloque I/P01/P02. Corregido y probado en CP1252 y UTF-8.\\
\bottomrule
\end{longtable}
\normalsize

\textbf{Regla de avance.} Cada actividad registra implementación, evidencia y
limitaciones antes de cerrarse. El orden puede ajustarse ante un defecto
bloqueante demostrado; cualquier cambio conserva IDs y justifica su dependencia.
Las pruebas y la documentación se actualizan en todos los bloques, no se aplazan
hasta P20. Las medidas de relevancia real necesitan modelos/proveedor y referencia
humana; si faltan, esa parte permanece pendiente y no se sustituye por un éxito
de pruebas simuladas.


\section{Ejecución del bloque I: P01--P04}
\label{sec:bloque1}
\textbf{Fecha: 18 de octubre de 2026.} Este bloque prepara una base verificable
antes de cambiar la interpretación de consultas o el refinamiento. No implementa
las rutas instrumentales, recuperación parcial o mejoras de score de bloques
posteriores. El registro siguiente complementa los cambios históricos 1--8.

\subsection{P01: perfiles y entorno reproducibles}
\textbf{Antes.} La lista de dependencias mezclaba ejecución, notebooks y
transitivas de una instalación no reproducida. El instalador mantenía otra
lista incompleta. La validación anterior dependía de paquetes instalados en
una carpeta temporal con metadatos superpuestos.

\textbf{Implementado.} Se crearon perfiles directos \code{core.in},
\code{models.in}, \code{indexing.in} y \code{dev.in} en \code{requirements/}.
Tras resolver e instalar en un entorno virtual vacío, se generaron cierres
exactos por plataforma: 55 paquetes en \code{core.lock}, 85 en
\code{models.lock} y 97 en \code{indexing.lock}. El perfil de desarrollo usa
el mismo runtime porque las pruebas utilizan unittest. El archivo principal
\code{requirements.txt} referencia el lock de inferencia, y el instalador
utiliza esos archivos y ejecuta \code{pip check}.

Se añadió \code{.env.example}, una referencia Python 3.12, configuración
centralizada y scripts de verificación, generación de locks y aceptación.
La instalación y sus perfiles se describen en \code{INSTALL.md}. Las rutas
de datos, prompt y salida del indexador se anclan al proyecto; el archivo
\code{.env} también se busca en esa carpeta.

\textbf{Evidencia.} Se verificaron 19 imports reales, incluidos Flask, FAISS,
LangGraph, Torch, Transformers, Sentence Transformers y el indexador. La carga
de los 258 registros y sus índices se validó sin descargar pesos o datasets.
Los informes están en \code{validation/block1/}. La reinstalación desde el
lock en un segundo entorno vacío completó 16 imports de inferencia, salud
correcta y \code{pip check} sin conflictos. Las 85 versiones del lock coinciden
exactamente con ese entorno; sus logs y manifiesto se conservan por separado.

\textbf{Por qué mejora.} Hace explícito qué se instala y permite comparar
versiones inequívocas. No se deduce que elegir versiones más nuevas mejore
la relevancia; se verifica compatibilidad de imports y control local.
\textbf{Límites.} El lock cubre Windows x64 y CPython 3.12. No fija revisiones
de pesos remotos ni hashes de todos los artefactos de distribución; no se
certifica otra plataforma o la inferencia con pesos reales.

\subsection{P02: fábrica Flask, contrato HTTP y configuración}
\textbf{Antes.} El endpoint podía operar sobre listas o tipos inválidos antes
de manejar errores; la importación WSGI no inicializaba el catálogo. El arranque
activaba debug sin configuración separada.

\textbf{Implementado.} \code{create_app} registra un blueprint e inicializa el
motor elegido. \code{wsgi.py} expone una instancia para Waitress; el arranque
local usa host, puerto y debug explícitos, con debug desactivado por defecto.
La fábrica permite inyectar un motor y desactivar la inicialización en pruebas.

\code{/search} comprueba tipo de contenido, JSON válido, objeto, texto no vacío
y longitud. Se devuelve 400 para entrada inválida, 413 para cuerpo excesivo,
503 para motor no listo o error controlado del flujo y 500 para errores
inesperados. El último caso registra el detalle en el servidor y entrega un
mensaje público genérico. Límites y agotamiento sin aprobación mantienen estado
HTTP 200 con advertencia y su causa de terminación.

\code{/health} devuelve 200 solo si el catálogo está listo y 503 si falla.
Esta salud acredita datos locales, no credenciales ni disponibilidad remota.
El servicio y la CLI siguen utilizando el mismo grafo.

\textbf{Evidencia.} Pruebas con Flask real cubren plantilla, salud, fábrica,
JSON malformado, tipos, longitud, cuerpo excesivo, búsqueda válida y estados
400/413/500/503. Se comprueba también arranque del catálogo real.
\textbf{Por qué mejora.} Evita errores de tipos no controlados y hace observable
la indisponibilidad. \textbf{Límites.} No se hizo una prueba de carga de servidor
ni una búsqueda completa contra Groq. Consumidores que asumían HTTP 200 ante
un error deben adaptarse al contrato documentado.

\subsection{P03: validez numérica y correspondencia del catálogo}
\textbf{Antes.} Las consultas podían quedar con norma nula tras sustraer un
negativo, y la carga no verificaba todos los valores o la coherencia del mapa.

\textbf{Implementado.} \code{vector_validation.py} comprueba forma, dimensión,
finitud y norma no degenerada de consultas, normalizando en precisión float64
y entregando un vector float32 contiguo. Los workers validan antes de llamar a
FAISS y conservan copias independientes de sus caches. Nomic conserva truncado
a 256 dimensiones; CLAP exige 512. Una consulta degenerada produce error
explícito, no un ranking sin significado.

Cada índice debe tener dimensión esperada, vectores presentes, estar entrenado
y usar producto interno. Se inspeccionan sus vectores por bloques: deben ser
finitos y unitarios con tolerancia de $10^{-3}$. El catálogo valida columnas,
IDs enteros no negativos y duplicados. El mapa debe cubrir todas las posiciones
y coincidir exactamente con el orden del catálogo; las cantidades de ambos
índices y CSV deben coincidir. La fábrica y los workers validan el conjunto
antes de intentar cargar los modelos.

\textbf{Evidencia.} Las pruebas usan índices FAISS reales y vectores con errores
de dimensión, métrica, valores NaN/infinito, norma nula y mapa inconsistente.
La consulta positiva idéntica a la negativa se rechaza antes de buscar.
Los dos índices locales pasaron las validaciones sobre sus 258 vectores.

\textbf{Por qué mejora.} Evita presentar datos incompatibles como recuperación
válida. \textbf{Límites y tareas derivadas.} La inspección añade coste de arranque;
es pequeña para este catálogo y debe medirse al crecer. El formato admite
reconstrucción de vectores y producto interno; otro tipo de índice necesita
revisar el validador. No se han resuelto todavía los presupuestos de P11 ni
la publicación atómica del catálogo de P17.

\subsection{P04: protocolo y línea base de control}
\textbf{Antes.} Existían pruebas de regresión y resultados históricos, pero no
una colección local fija para comparar las siguientes modificaciones de forma
reproducible con hashes de código y datos.

\textbf{Implementado.} \code{evaluation/queries.json} fija 12 consultas, dos por
categoría: letras, audio, mixtas, literales, negativas e instrumentales.
El runner permite baseline sin reescritura y flujo con refinamiento, con modos
simulado y real explícitos. Registra versiones, hashes, planes, rankings,
estados, exclusiones, trazas, latencia y contadores por etapa. Los fixtures
de control generan rankings y veredictos deterministas; no son juicios musicales.

Se realizaron tres repeticiones por consulta: 36 corridas sin refinamiento y
36 con refinamiento simulado, ambas con cero errores, duplicados y reapariciones
de excluidos. La estabilidad Jaccard fue 1 en este fixture determinista; no
demuestra estabilidad de un LLM real ni calidad de orden. Las métricas de
relevancia y coste monetario permanecen nulas, no ceros inventados.

\code{qrels.json} declara juicios humanos pendientes. El protocolo rechaza IDs
desconocidos, grados inválidos y cobertura completa declarada sin anotar todo
el catálogo. Recall/nDCG requieren cobertura suficiente; la tasa de violaciones
requiere etiquetas explícitas. La simulación nunca usa las etiquetas para
afirmar calidad musical. El modo real se registró como bloqueado por falta
de clave Groq; no se hicieron llamadas externas.

\textbf{Por qué mejora.} Hace comparables las futuras modificaciones y evita
confundir regresiones de control con evaluación musical. \textbf{Límites.}
La línea base corresponde al estado posterior a P01--P03 y anterior a P05,
no al código original A/B/C. Faltan anotación humana, pesos y mediciones reales;
esa parte permanece pendiente para P12/P20 y no se da por realizada.
Los contadores de etapas no incluyen reintentos internos del proveedor y no
permiten inferir coste monetario.

\subsection{Estado del bloque y tareas que permanecen abiertas}
P01--P03 tienen implementación y verificación local. P04 tiene protocolo y
línea base simulada; la evidencia musical real sigue pendiente por los
requisitos anteriores. No se atribuye mejor relevancia o latencia de inferencia.
E11 queda cubierto para el entorno y HTTP verificados, manteniendo pendiente
la validación integral de P20. E10 queda cubierto en su parte numérica;
presupuestos y exploración siguen en P11. E07 conserva pendiente la comparación
real del refinador en P12.

La siguiente implementación sigue siendo P05--P08. Las limitaciones anteriores
se conservan con sus IDs; no se cierran por disponer ahora de más pruebas.


\section{Ejecución del bloque II: P05--P08}
\label{sec:bloque2}
\textbf{Fecha: 18 de octubre de 2026.} Se implementaron propiedad del estado,
rutas por modalidad, procedencia de restricciones y un fallback conservador.
Se mantienen pendientes las mediciones con modelos reales y anotaciones humanas.
Este bloque cambia algunos estados de la API y reduce deliberadamente la
cobertura automática de interpretaciones inciertas; no es una mejora universal.

\subsection{P05: estado serializable y propietario único de exclusiones}
\textbf{Antes.} Juez y grafo podían escribir el pool. La idempotencia evitaba
duplicados, pero el estado contenía un objeto de pool y una caché, inadecuados
como contrato de persistencia.

\textbf{Ahora.} El juez devuelve decisiones sin mutar exclusiones. El único
punto de escritura del flujo es el nodo de evaluación del orquestador.
El estado contiene listas y diccionarios serializables: decisiones, exclusiones,
versión de plan, candidatos y evidencia. La caché de recuperación queda fuera
del estado, ligada a la instancia de grafo creada para una solicitud.
El antiguo argumento de pool se admite por compatibilidad, pero se ignora;
el tipo auxiliar no participa en el flujo actual.

Cada exclusión registra candidato, motivo, iteración, alcance y versión del
plan. Los datos inexistentes se excluyen para toda la solicitud; las decisiones
dependientes de las consultas se vinculan al plan evaluado. Un rechazo histórico
no reaparece como alternativa sin evaluar: solo el candidato seleccionado puede
volver después de reevaluarse con un plan nuevo. Se conserva la decisión previa
en el historial.

\textbf{Evidencia.} Pruebas de estado convertible a JSON sin objetos opacos,
decisión pura del juez, escritura única por rechazo y solicitudes concurrentes
con estados separados. \textbf{Ventaja.} Se puede auditar quién decidió qué y
con qué plan; desaparece la doble responsabilidad de mutación.
\textbf{Límite.} No se añadió almacenamiento durable ni reanudación. La caché
fuera del estado no se guarda al serializar. Reutilizar el grafo o añadir
checkpoints necesita el trabajo de P16/P18.

\subsection{P06: modalidades e instrumentales sin aprobación acústica ficticia}
\textbf{Antes.} El flujo exigía letras a cualquier candidato e invocaba al juez
incluso para consultas acústicas. Una aprobación lírica podía confundirse con
una aprobación multimodal.

\textbf{Ahora.} Una consulta solo acústica omite el juez lírico, admite candidatos
sin letra y termina como \code{retrieved}: recuperado, con audio
\code{not_evaluated}. Una consulta solo lírica puede terminar como
\code{approved}. Si la letra se aprueba y también se pidió audio, la terminación
es \code{partially_validated}; el sonido sigue sin evaluación independiente.

Cada resultado informa \code{validation.lyrics} y \code{validation.audio}, con
estados de aprobado, rechazado, error, no solicitado o no evaluado. La interfaz
muestra esos estados. Una ausencia de letra en una solicitud que sí la requiere
se registra como falta de evidencia, no como rechazo de un LLM que no fue llamado.
\code{judge_iterations} cuenta invocaciones reales al juez; el contador de
selecciones queda en \code{candidate_iterations}.

\textbf{Evidencia.} Instrumentales sin letra, omisión de llamadas al juez,
validación parcial en consultas mixtas y respuestas de Flask real.
\textbf{Ventaja.} Evita descartar instrumentales por una condición ajena a su
modalidad y declara el alcance de la evidencia.
\textbf{Límite.} CLAP sigue proporcionando recuperación, no certificación de
instrumentos, voces o energía. La calidad de la ruta acústica no se ha medido.
Las restricciones líricas negativas también requieren evaluación aunque no haya un tema positivo lírico. El presupuesto todavía cuenta selecciones; su separación completa sigue en P11.

\subsection{P07: restricciones con origen y reparación del plan}
\textbf{Antes.} Todos los negativos quedaban inmóviles, aunque fueran una
inferencia equivocada. Tampoco podía activarse una modalidad omitida por el router.

\textbf{Ahora.} \code{intent.py} prepara restricciones con ID, modalidad,
polaridad, texto, origen y cita. \code{user_explicit} exige una cita literal del
usuario; una negativa debe incluir una negación reconocida. Lo no fundamentado
se registra como \code{model_hypothesis}. Las restricciones explícitas capturadas
se conservan; las hipótesis pueden corregirse o retirarse.

Una modalidad puede activarse o desactivarse para corregir la intención si se
aporta una cita y diagnóstico de esa corrección. Cambiar audio siempre requiere
evidencia textual con una señal de modalidad soportada: un rechazo lírico por
sí solo no justifica describir como incorrecto el sonido de una canción no oída.
Un cambio de plan se versiona y sus decisiones se reevaluarán bajo esa versión.
Las modalidades no seleccionadas permanecen intactas.

Se rechaza un plan que omita anclar una negación visible del usuario. Este
control usa señales léxicas, no una comprensión completa del lenguaje. El prompt
del router ahora pide procedencia y listas vacías para modalidades no solicitadas.
Se probó el formateo con la plantilla LangChain real para que los ejemplos JSON
no introduzcan variables de plantilla inexistentes.

\textbf{Evidencia.} Pruebas de eliminación de hipótesis, conservación de
restricciones explícitas, citas inventadas o sin negación, corrección de
modalidades y reevaluación después de cambiar de plan.
\textbf{Ventaja.} Separa decisiones revisables de restricciones que el usuario
declaró y evita perderlas silenciosamente.
\textbf{Límite importante.} Verificar una cita no demuestra que su traducción
inglesa sea semánticamente fiel. Un modelo todavía puede inferir demasiado al
etiquetar una restricción como explícita. Se registra E13 para calibrarlo.

\subsection{P08: fallback limitado, trazable y con abstención}
\textbf{Antes.} El respaldo copiaba la consulta cruda en ambas modalidades,
con negativos vacíos, sin garantía de idioma o intención.

\textbf{Ahora.} Puede reutilizarse un plan válido únicamente si está asociado
a la misma consulta exacta. La función permite esa recuperación, pero no guarda
un historial global de consultas ni se conecta todavía a una persistencia.
Cuando no hay plan previo, una gramática local acotada interpreta frases de
sonido soportadas en español e inglés, conserva las negaciones reconocidas y
registra sus citas. Por ejemplo, piano lento sin voces produce consultas
acústicas positivas y una negativa de voces, con la rama lírica vacía.

Si quedan palabras o alcances de negación sin interpretar, se devuelve
\code{clarification_needed}, listas de búsqueda vacías y una advertencia de
reformulación. No se llama a workers ni al juez. No se promete traducción
general, interpretación lírica offline ni extracción arbitraria de negativos.
La respuesta HTTP es 200 con ese estado: el JSON es válido, pero no se ejecutó
una búsqueda de significado incierto.

\textbf{Evidencia.} Frases bilingües soportadas, restricciones negativas,
negaciones ambiguas, planes previos de otra consulta y fallo del router.
\textbf{Ventaja.} La degradación no presenta como fiel una consulta que perdió
restricciones. \textbf{Coste.} Hay consultas para las que el respaldo devuelve
menos resultados o ninguno y exige reformular; la cobertura y satisfacción
necesitan medición antes de ampliar reglas.

\subsection{Resultados de control y regresiones observadas}
La suite pasó 72 pruebas en el entorno del bloque I, con Flask, FAISS y
LangGraph reales y proveedores/modelos simulados. Se actualizaron expectativas
anteriores para los estados nuevos y el juez sin efectos laterales; no se
eliminaron los casos de regresión originales.

Se ejecutaron 36 corridas simuladas sin refinamiento y 36 con refinamiento.
En cada serie se observaron 12 aprobaciones líricas, 6 validaciones parciales,
12 recuperaciones acústicas sin validación independiente y 6 solicitudes de
reformulación. Estas últimas corresponden a dos consultas repetidas tres veces:
\code{negative_romance} y \code{instrumental_study}. Sus planes sintéticos
anteriores no incluyen la procedencia requerida para las negaciones y la
gramática limitada no cubre toda la petición. El fixture original se mantuvo
para que ese cambio de cobertura sea visible.

No hubo errores del flujo. Se registran por separado reconsideraciones históricas
y reapariciones de excluidos sin reevaluación; las segundas deben permanecer en
cero. Jaccard solo se calcula cuando hay resultados; una respuesta sin ranking
no recibe un 1 ficticio. Las métricas musicales continúan nulas. La ejecución
real se mantuvo bloqueada sin llamar al proveedor.

Esto demuestra controles adicionales, no mejor relevancia musical. La reducción
de cobertura automática es una contrapartida explícita; E02/P12 deben evaluar
si conviene ampliar el respaldo o cambiar la estrategia con evidencia real.
Los informes del bloque I se conservaron; los actuales están en
\code{validation/block2/} y llevan el identificador \code{block2-pre-P09}.

\subsection{Compatibilidad y trabajo que no se cierra}
La API añade \code{retrieved}, \code{partially_validated} y
\code{clarification_needed}; clientes que esperaban solo aprobación/error/límite
deben interpretar estos estados. \code{validation_scope} puede ser
\code{none}; \code{constraints}, \code{plan_version} y validación por resultado
permiten inspeccionar la evidencia. El argumento legado de pool ya no aplica
exclusiones por sí mismo. Los locks no cambiaron y no se instalaron dependencias.

P05--P08 tienen implementación y verificación de control. Se cubren la propiedad
única de E08, las rutas de E01 y los controles de E02/E03. Persistencia, evidencia
acústica independiente, fidelidad semántica y evaluación real siguen abiertas.
Los siguientes puntos de implementación son P09--P12, no la ampliación de la UI
o ingesta. La recuperación parcial ante fallos del refinador sigue pendiente.

\section{Ejecución del bloque III: P09--P12 y correcciones asignadas}
\label{sec:bloque3}
\textbf{Fecha: 18 de octubre de 2026.} Se implementó el control local de fallos,
evidencia y presupuestos. P10 permanece abierto para calibración humana y P12
para medir utilidad musical real. No se afirma superioridad de relevancia.

\subsection{Reubicación y corrección de E14--E16}
E14 pertenece al bloque II/P07: la identidad del plan debe distinguir cambios
efectivos de una permutación de listas. La huella ahora canoniza descriptores y
registros de restricciones sin borrar procedencia. La caché de recuperación usa
la misma independencia de orden. Una prueba del grafo rechaza candidatos 1 y 2
una sola vez, conserva la versión 1 y termina por ventana agotada; antes repetía
el candidato 1 tres veces y creaba versiones 2 y 3. Un descriptor efectivamente
distinto sigue cambiando la identidad. Esto no demuestra equivalencia semántica
entre textos distintos: E13 continúa pendiente.

E15 pertenece al bloque II/P05/P08: la salida de aclaración ahora verifica JSON
finito igual que un plan listo. Números NaN/infinito y objetos opacos se rechazan;
el orquestador conserva su fallback controlado. Se probaron ambas rutas y la
serialización del estado final recuperado. Se conserva el rechazo explícito del
dato inválido, sin sustituirlo por una confianza inventada.

E16 pertenece al bloque I/P01/P02: los mensajes de router, workers y juez usan
una salida que escapa símbolos incompatibles con la codificación de consola.
Un diagnóstico deja de interrumpir una inicialización correcta. UTF-8 sigue
siendo una opción operativa; en CP1252 algunos símbolos se ven escapados.
Ambas codificaciones completaron las mismas 96 pruebas.

\subsection{P09: recuperación parcial y continuidad}
\textbf{Antes.} Una excepción de una rama abortaba la recuperación completa;
un refinador inválido detenía el flujo. \textbf{Ahora.} Cada rama informa éxito,
vacío, ausencia de solicitud o fallo. Los resultados útiles de la otra rama se
conservan; sus pesos efectivos se renormalizan y se exponen junto a los pesos
solicitados y la política. Las restricciones originales no se eliminan.
Si falla el refinador, se conserva el plan válido y se deshabilitan nuevos
intentos del refinador para esa solicitud. Si falla toda una nueva recuperación,
se restaura el plan, versión y ranking obtenidos juntos en la última búsqueda
válida, aplicando las exclusiones de esa versión.

\textbf{Evidencia y contrapartida.} Pruebas del grafo y HTTP cubren fallos
parciales, fallos de ambas ramas y restauración. Una petición mixta sigue
informando sonido sin evaluar aunque la letra se apruebe. Esta política mejora
continuidad, pero puede reducir cobertura y cambiar scores; no garantiza la
misma calidad que disponer de ambas ramas. P13 debe explicar contribuciones.

\subsection{P10: evidencia lírica, abstención y reintento acotado}
\textbf{Antes.} El juez veía los primeros 3000 caracteres y emitía texto o una
herramienta de rechazo. \textbf{Ahora.} Se solicita JSON con estado, motivo y
citas literales. Las citas se contrastan con fragmentos realmente suministrados.
Las letras de hasta 12000 caracteres se incluyen completas; las más largas
aportan comienzo, centro y final dentro de ese presupuesto. La cobertura parcial
impide aprobar el conjunto. La abstención no se transforma en rechazo semántico
ni activa el refinador; puede continuar la revisión de otros candidatos.

Un formato o evidencia inválidos permiten un solo reintento correctivo;
una excepción del proveedor no se reintenta por esta capa. El cliente del juez
desactiva sus reintentos internos y registra intentos de modelo. Se conserva
compatibilidad con veredictos antiguos, etiquetados como \texttt{legacy}, sin
atribuirles citas verificadas. Esta fue la decisión inicial del bloque;
la sección~\ref{sec:correcciones_pendientes} documenta su retirada posterior.
Los demás candidatos siguen sin evaluar.

\textbf{Evidencia y contrapartida.} Una prueba encuentra y cita evidencia
después del carácter 4000; otras cubren citas inventadas, cobertura incompleta,
abstención, errores y el máximo de dos invocaciones. El contexto puede ser mayor
que antes y el máximo de salida pasó de 256 a 1024 tokens; ello puede aumentar
coste y latencia. Una cita exacta no prueba que el juicio semántico sea correcto.
\texttt{evaluation/judge\_metrics.py} calcula falsos aprobados/rechazados,
abstención y error solo con etiquetas humanas. Sin pares anotados devuelve
métricas nulas: la calibración y E13 siguen pendientes, no se inventaron etiquetas.

\subsection{P11: presupuestos independientes}
\textbf{Antes.} Tres visitas a candidatos podían agotarse sin llamar al juez
si faltaban metadatos; se buscaba una ventana fija. \textbf{Ahora.} Por solicitud
hay límites separados: 20 visitas, 3 evaluaciones del juez (hasta 6 intentos de
modelo por el reintento), 2 refinamientos y 1 ampliación de recuperación de
top-10 a top-20 por rama. Las omisiones de metadatos no consumen evaluaciones
del juez. Los contadores y el motivo de parada se exponen en la API.

La ampliación solo se intenta al agotar una ventana saturada; se filtran las
exclusiones previas. Agotar esa ventana no demuestra ausencia de canciones
válidas en todo el catálogo; la API no declara agotamiento global.
\textbf{Pruebas.} Cuatro metadatos ausentes antes de un candidato válido,
presupuesto total de visitas, ampliación que encuentra un candidato nuevo y
agotamiento de la ampliación. Revisar más candidatos puede consumir más trabajo;
los límites son configurables al construir el grafo y no una estimación monetaria
ni un límite global de todas las solicitudes HTTP del proveedor.

\subsection{P12: comparación pareada y política de refinamiento}
\textbf{Antes.} El refinador se invocaba tras rechazos sin evidencia de mejora
musical. \textbf{Ahora.} Está desactivado por defecto en la aplicación; la opción
\texttt{SPELLFIND\_ENABLE\_REFINEMENT=true} permite uso experimental. Solo se
considera tras rechazo lírico, con contrato válido y presupuesto disponible.
No se activa por abstención, metadatos ausentes ni recuperación acústica sola.

Se capturaron 36 ejecuciones sin refinamiento y 36 con él, pareadas por consulta
y repetición, con hashes iguales de fuentes, tests, datos y locks. Ambas series
simuladas tuvieron 12 aprobaciones líricas, 6 validaciones parciales,
12 recuperaciones acústicas y 6 aclaraciones, sin errores. El tratamiento añadió
18 invocaciones al refinador y 18 búsquedas líricas; no cambió el número de
evaluaciones del juez. Los veredictos y rankings son sintéticos: no demuestran
mejor o peor relevancia. Los tiempos simulados y su dispersión se conservan en
el informe, sin convertirlos en latencia de modelos o precio del proveedor.

\texttt{evaluation/compare.py} rechaza snapshots distintos y muestra diferencias
pareadas. La evaluación real escribió estado bloqueado por falta de
\texttt{GROQ\_API\_KEY}; las etiquetas humanas siguen vacías. Mantener la opción
desactivada limita trabajo adicional no justificado, pero renuncia a posibles
beneficios todavía no medidos. Para decidir su activación rutinaria se requieren
comparaciones reales de relevancia, violaciones negativas, abstención y coste,
con anotaciones y condiciones comparables. P12 no queda cerrado por la simulación.

\subsection{Verificación, compatibilidad y pendientes}
\textbf{Archivos.} E14/E15 cambian \texttt{intent.py} y la clave de caché en
\texttt{app.py}; E16 añade \texttt{console.py} y lo integra en router, workers
y juez. P09/P11 usan \texttt{workflow.py}, \texttt{retrieval.py} y \texttt{app.py}.
P10 cambia \texttt{judge.py}; P12 actualiza aplicación, ejemplo de configuración
y \texttt{evaluation/run.py}, y añade comparación y métricas de calibración.
La interfaz reconoce abstención. \texttt{tests/test\_block3.py} añade regresiones;
se actualizan expectativas anteriores para la continuidad y presupuestos nuevos.

\texttt{scripts/verify\_block3.py} ejecutó 96 pruebas con salida CP1252 y 96 con
UTF-8, ambas series correctas. La evidencia está en \texttt{validation/block3/}.
No se instalaron dependencias, descargaron pesos ni hicieron llamadas al modelo.
Los informes de bloques anteriores y el diagnóstico original conservan su
carácter histórico. El documento se actualizó en fuente; no se certifica en esta
entrega una compilación PDF.

Los estados y campos añadidos requieren que clientes reconozcan abstenciones,
degradaciones y presupuestos. El nombre histórico \texttt{judge\_iterations}
sigue contando evaluaciones; los intentos de modelo están en contadores separados.
P13--P16 son el siguiente bloque de implementación. P10/P12 y la fidelidad
semántica E13 mantienen sus verificaciones abiertas; no se renumeran ni se
consideran resueltas al avanzar.

\section{Ejecución del bloque IV: P13--P16}
\label{sec:bloque4}
\textbf{Fecha: 18 de octubre de 2026.} P13--P15 están implementados y probados
localmente. P16 incorpora instrumentación y un perfil del runtime con catálogo
e índices reales y modelos de control; sus mediciones con modelos/proveedor
reales y memoria completa del proceso siguen pendientes. P04/P10/P12 y E13
no se consideran cerrados por avanzar en este bloque.

\subsection{P13: puntuación explicable y pesos configurables}
\textbf{Antes.} Se combinaban scores de workers normalizados por mínimo/máximo
con una bonificación RRF independiente de los pesos. El resultado se multiplicaba
por cien y podía superar cien, sin ser una probabilidad. Una rama con peso cero
podía seguir aportando por la bonificación.

\textbf{Ahora.} \texttt{scoring.py} define una fusión por posición:
\[
R_m(d)=\begin{cases}\frac{61}{60+r_m(d)},&d\text{ aparece en la rama }m,\\
0,&\text{en otro caso},\end{cases}
\qquad S(d)=w_L R_L(d)+w_A R_A(d).
\]
La posición comienza en uno. Los pesos efectivos suman uno entre ramas con peso
positivo que aportan resultados. Se exponen rango, score bruto del worker,
score por posición, peso y contribución por canción; las contribuciones suman
el score final. Los empates se resuelven por ID, de forma determinista.

La API versión 2 elimina \texttt{score\_pct} y añade el método
\texttt{weighted\_normalized\_rrf\_v2}, rango $[0,1]$ e indicador de que el score
no es probabilidad. La interfaz presenta decimales y explica el significado.
Las decisiones del juez pueden priorizar una recomendación frente a otra con
mayor score: se declara esa política y su motivo por resultado.

Los pesos opcionales de solicitud son dos números finitos entre cero y uno que
suman uno. Se distinguen pesos solicitados y efectivos, con ajuste explícito
ante ramas ausentes, vacías o fallidas. Peso cero desactiva la recuperación de
esa rama, sin eliminar conceptos o restricciones que debe comprobar el juez.
No se inventa una modalidad cuando el plan carece de consultas para ella.

\textbf{Ventaja y contrapartida.} Se elimina el porcentaje engañoso y el aporte
de ramas con peso cero. La nueva fusión evita comparar magnitudes de espacios
distintos, pero descarta diferencias de magnitud dentro de cada ranking y puede
cambiar el orden respecto a la versión anterior. Los primeros puestos tienden
a valores cercanos a uno; ello no demuestra relevancia alta. Scores de versiones,
consultas y ventanas distintas no se comparan como probabilidades. La superioridad
musical de esta fórmula sigue pendiente de P04/P12; no se infiere de sus contratos.

\subsection{P14: filtros duros por metadatos}
\textbf{Antes.} No había selección exacta por artista, género o fechas.
\textbf{Ahora.} \texttt{search\_options.py} valida artista y género completos,
ignorando mayúsculas y espacios exteriores y normalizando Unicode NFC, sin
eliminar acentos ni hacer coincidencias parciales. El género procede de
\texttt{tag}; no se infiere del audio. Los años son enteros y los límites son
inclusivos. Los filtros se combinan con AND; años ausentes o inválidos no
satisfacen un filtro temporal. Columnas opcionales ausentes no rompen los
controles de catálogo.

\texttt{GET /catalog/options} ofrece nombres y límites de años sin letras ni
llamadas a modelos. Con filtros activos se recupera hasta el tamaño del catálogo
antes de filtrar cada rama y fusionar la ventana elegible. Así, una canción en
posición 25 del ranking global puede ser la primera coincidencia filtrada.
El juez solo recibe candidatos elegibles. Cero coincidencias de metadatos se
resuelve antes del router, sin consumir modelos, y queda distinguido de una
consulta semántica sin resultados utilizables.

\textbf{Evidencia y límites.} Pruebas HTTP cubren esa canción fuera del top-10,
intersección, tipos inválidos, límites inclusivos y ausencia de coincidencias.
La operación usa un catálogo validado y estable durante el proceso; P17 sigue
siendo necesario para actualización coordinada. Recuperar todos los candidatos
es sencillo para 258 canciones, pero aumenta trabajo y memoria al crecer el
catálogo. Evaluar preselección nativa en el índice o índices de metadatos dentro
de P16/P17 antes de escalar. La completitud de metadatos tampoco certifica su
exactitud ni su género acústico.

\subsection{P15: evidencia visible y alcance por candidato}
\textbf{Antes.} La respuesta descartaba letras completas y la tarjeta mostraba
principalmente título, score y un estado, sin fragmentos ni contribuciones.
\textbf{Ahora.} \texttt{result\_evidence.py} selecciona hasta 320 caracteres.
Prioriza citas estructuradas que ocurren literalmente en la letra. Sin esas
citas, usa una heurística léxica de las consultas positivas o un fragmento
inicial, rotulados como orientativos. Conserva offsets del texto original y
advierte recortes. Las solicitudes puramente acústicas no reciben una supuesta
prueba lírica. No se envía la letra completa como parte de cada resultado.

Cada tarjeta muestra contribuciones, metadatos, alcance de letra/audio, motivo
del juez cuando existe y distinción entre recomendación evaluada y alternativa
recuperada. Se conserva el estado de error de un candidato realmente evaluado;
no se reemplaza por un falso estado de candidato nunca examinado. Se usa solo
evidencia de la versión de plan correspondiente. Las decisiones heredadas sin
citas mantienen esa limitación explícita. El contenido variable se inserta como
texto DOM, incluyendo títulos, citas y mensajes.

\textbf{Evidencia y límites.} Pruebas cubren citas reales/inventadas, offsets
Unicode, ausencia de aprobación heredada en alternativas y contexto acústico.
La web se verificó en un navegador con fixture visiblemente marcado como
control: catálogo de 258 canciones, rankings y juez simulados, filtros
combinados, pesos personalizados, año inválido y cero coincidencias. La captura
\texttt{validation/block4/ui-desktop.jpg} permite revisar el resultado. Hay CSS
adaptativo; no se certifica una revisión visual móvil. Una coincidencia textual
o un fragmento orientativo no prueban cumplimiento semántico: E13/P10 siguen abiertos.

\subsection{P16: instrumentación y decisiones basadas en mediciones}
\textbf{Antes.} No había desglose por construcción del grafo, ejecución,
presentación, codificación, índice y contención. \textbf{Ahora.} La API informa
tiempos del servicio y nodos. Las rondas de recuperación registran caché de la
solicitud, perfil de ejecución de workers y fallos. En un acierto de caché no se
atribuyen otra vez los tiempos del worker anterior. Los workers miden acceso a
recursos, espera del bloqueo, codificación/caché, búsqueda FAISS y RRF interno.
Los perfiles usan contexto de ejecución separado para solicitudes concurrentes.
Los tiempos internos no incluyen necesariamente serialización HTTP ni red.

\texttt{evaluation/profile.py} ejecuta una primera consulta, 20 consultas
calientes secuenciales y 20 con dos y con cuatro hilos. Usa los índices FAISS
y metadatos reales; los encoders reconstruyen vectores existentes y router/juez
son controles. Por tanto, se mide el runtime local, no inferencia Nomic/CLAP ni
servicio Groq. Se fija un hilo OpenMP por tarea solo en el experimento. Reporta
medias, mediana, p95, dispersión, tasa de solicitudes, contención y cachés.
Cinco solicitudes adicionales miden asignaciones Python con
\texttt{tracemalloc} en una fase separada, para no sesgar los tiempos calientes
frente a la primera consulta sin trazado. El total es 66 solicitudes de control.

El experimento verificó identidad de resultados entre las condiciones y
reutilización de cachés. La memoria trazada excluye tensores nativos, FAISS y
GPU, y no equivale a RSS; se informa separadamente el payload de vectores del
catálogo. Las métricas de memoria completa y modelos reales permanecen nulas.
Los valores y hashes reproducibles están en \texttt{validation/block4/profile.json}.

\textbf{Decisión y pendiente.} Se mantienen grafos por solicitud y cachés
existentes acotadas; no se introduce un grafo compartido, otra caché global ni
batching por una aceleración aparente en modelos de control. Un grafo reutilizado
exigiría separar su caché de cada solicitud y medir invalidación/contención.
Comparar con máquina de estados simple solo si la compilación pesa realmente en
la carga final. P16 continúa abierto para carga con modelos, memoria completa y
optimización demostrada; E04/E09 y P17/P18 conservan su trazabilidad.

\subsection{Archivos, compatibilidad y verificación}
\texttt{app.py} integra score, filtros, evidencia y mediciones; \texttt{cli.py}
admite peso lírico, artista, género y años. Se añaden
\texttt{scoring.py}, \texttt{search\_options.py}, \texttt{result\_evidence.py} y
\texttt{profiling.py}. Se instrumentan \texttt{workflow.py} y ambos workers.
La plantilla contiene controles y tarjetas actualizadas. Las pruebas nuevas
están en \texttt{tests/test\_block4.py}; los informes y capturas, en
\texttt{validation/block4/}. \texttt{scripts/preview\_block4.py} es exclusivamente
un fixture local de interfaz, señalado como simulación; no es configuración de
producción ni evaluación musical.

Se ejecutaron 115 pruebas en CP1252 y 115 en UTF-8. Se mantienen capturas de
36 simulaciones sin refinamiento y 36 con él, sin modificar las consultas
históricas ni sustituir juicios humanos. La evaluación real sigue bloqueada
cuando falta la clave. El perfil adicional no descarga pesos ni llama al
proveedor. No hay dependencias nuevas. Las fuentes TeX se verifican por
inclusiones y referencias; no se certifica una nueva compilación PDF.

Los clientes deben migrar de \texttt{score\_pct} a \texttt{score} y sus
contribuciones; el endpoint sigue admitiendo solicitudes antiguas con solo
\texttt{query}. La firma de consola comparte los mismos contratos de opciones.
Los informes anteriores conservan fórmulas y snapshots de su fecha: no comparar
scores numéricos entre métodos como si compartieran escala semántica.

El siguiente bloque de implementación es P17--P20. Permanecen pendientes las
validaciones reales de P04/P10/P12/P16 y la fidelidad E13. El bloque IV aporta
transparencia y controles, sin atribuirse relevancia o rendimiento de modelos
todavía no medidos.

\section{Bloque V: catálogo actualizable y operación integrada}
\label{sec:bloque5}
\textbf{Fecha: 18 de octubre de 2026.} Implementación local de P17--P20.
La evidencia se conserva en \code{validation/block5/}. Las pruebas con adaptadores
de control verifican contratos y fallos, no relevancia musical ni descargas reales.

\subsection{P17: generaciones coherentes del catálogo}
Antes se conservaban caches independientes del CSV, mapa e índices. Reemplazar
archivos mientras una búsqueda estaba en curso podía combinar generaciones.
Ahora \code{catalog.py} publica los cuatro artefactos en un directorio inmutable,
valida dimensiones, alineación, normalización y hashes SHA-256 y cambia el puntero
actual mediante una transacción SQLite. El historial publicado distingue las
generaciones válidas de los directorios huérfanos de una publicación fallida.
El primer arranque copia los archivos heredados de \code{output/}, sin modificarlos.

Cada solicitud web o búsqueda CLI fija un snapshot durante todo su recorrido.
Una consulta iniciada antes de una actualización termina con su generación; la
siguiente adopta la nueva. Los modelos se reutilizan, pero los recursos de búsqueda
se superponen desde el contexto de la solicitud. No se agregan caches globales de
resultados. Los helpers heredados sin contexto conservan su comportamiento estático.
Las vistas de metadatos son copias y los índices expuestos no permiten agregar ni
borrar vectores. Ante un nuevo puntero inválido, un servidor que ya posee un snapshot
continúa con él e informa una advertencia; un arranque sin catálogo válido falla.

Los publicadores se serializan: agregar canciones recarga la última generación
dentro de la transacción. Reconstruir todo exige que la versión inicial no haya
cambiado; así una ingesta concurrente no desaparece silenciosamente. Los enlaces
de reproducción incluyen versión e ID, evitando que un ID reutilizado reproduzca
otra canción. Se comprueba que la versión pertenezca al historial publicado.
El coste es copiar índices y catálogo en cada publicación: es apropiado para el
catálogo pequeño actual, pero debe medirse antes de escalar. No se garantiza
durabilidad frente a pérdida de energía ni se implementa todavía rollback operativo.

\subsection{P18: decisión sobre persistencia y progreso}
No se recuperan checkpoints del grafo de búsqueda: sus juicios y recursos podrían
pertenecer a otra generación, y no existe evidencia de que el ahorro compense esa
complejidad. La búsqueda sigue siendo independiente por solicitud. Para ingesta,
sí se persisten trabajos en SQLite: estado, etapa, progreso, metadatos, propietario,
lease, resultado y error. La interfaz consulta el estado cada 1,5 segundos y conserva
el último ID en almacenamiento local. Recargar consulta ese ID; no vuelve a enviar
la canción ni reanuda automáticamente la búsqueda.

Un trabajo pendiente puede ser reclamado una sola vez entre procesos. Un trabajo
en ejecución pierde su propiedad al vencer el lease de 60 segundos y queda
interrumpido. Un trabajador vencido no puede renovar ni publicar. El usuario puede
enviar de nuevo el enlace; la publicación es idempotente por trabajo y URL. Si hubo
una caída tras publicar y antes de registrar el resultado, el nuevo envío detecta
la canción existente. La cola admite hasta 32 trabajos activos y deduplica enlaces.
Se eligió polling por sencillez y recuperación por ID; SSE queda diferido hasta
demostrar una necesidad de menor latencia o menor tráfico.

\subsection{P19: ingesta opcional sin bloquear la búsqueda}
La función permanece desactivada por defecto mediante
\code{SPELLFIND_INGEST_ENABLED=false}. El adaptador nativo usa yt-dlp para MP3 con
FFmpeg, prioriza subtítulos oficiales y utiliza Whisper CPU como alternativa.
Un fallo al descargar subtítulos no descarta el audio ya descargado. Se aceptan
enlaces HTTPS de vídeos individuales de YouTube, con duración conocida de hasta
15 minutos, límites de tamaño y tiempos de espera. Se rechazan directos y enlaces
arbitrarios. Estos límites no equivalen a cuotas globales de almacenamiento.

Las letras se codifican con prefijo \code{search_document:}, fragmentos de 2000
caracteres y media normalizada de vectores de dimensión 256. CLAP procesa ventanas
de 10 segundos al inicio, centro y final a 48 kHz y produce una media normalizada
de dimensión 512. La cobertura es parcial y se registra. La receta del catálogo
importado no demuestra igualdad de revisiones ni de agregación con los vectores
nuevos; esa compatibilidad debe evaluarse antes de afirmar calidad equivalente.

Sólo se publica después de completar metadatos y ambos vectores. No se inventan
letras vacías ni se usa la fecha de subida como año de lanzamiento. Artista, género
y año pueden ser aportados explícitamente; la procedencia queda registrada. Una
canción sin transcripción utilizable falla, incluso si es instrumental: el esquema
alineado actual requiere ambas modalidades. Esto es una limitación pendiente,
no una mejora sobre recuperación por audio de instrumentales.

La API añade \code{GET /ingest/status}, \code{POST /ingest} y
\code{GET /ingest/jobs/<id>}. Responde 202 al encolar, 400 ante datos inválidos,
403 ante origen cruzado, 429 ante saturación y 503 si la capacidad está desactivada
o incompleta. La interfaz muestra etapas, errores y reproducción versionada.
El servicio está concebido para uso local; la comprobación de origen no sustituye
autenticación para exposición pública. La extensión
\code{requirements/ingestion.in} es experimental y no tiene un lock verificado.
En este entorno faltan yt-dlp, FFmpeg y Whisper: no se realizó una ingesta real.
Referencias técnicas: \url{https://github.com/yt-dlp/yt-dlp},
\url{https://github.com/openai/whisper} y \url{https://github.com/LAION-AI/CLAP}.

\subsection{P20: comprobaciones y límites del cierre}
Las regresiones ejercitan publicación fallida, hashes corruptos, snapshots
inmutables, dos publicadores, comparación de versión, búsqueda en curso durante
actualización, propiedad y vencimiento de trabajos, errores del adaptador,
subtítulos y fallback, contratos HTTP, audio con rangos y consistencia web/CLI.
Las evaluaciones fijan un snapshot para toda la serie y registran su manifest;
la comparación rechaza catálogos diferentes. Los resultados ejecutados y sus
comandos quedan en los informes del bloque. No se confunden adaptadores de
control con ejecución de proveedores ni una comprobación estática con compilación PDF.

\subsection{Pendientes adicionales y prioridad}
\begin{description}
\item[E17, P17/P20: retención y recuperación.] Definir cuotas, limpieza segura de
medios, trabajos y generaciones huérfanas, backups y restauración. Actualmente se
retienen para no eliminar recursos de consultas o enlaces todavía activos.
\item[E18, P01/P19/P20: entorno e integración reales.] Fijar dependencias opcionales
en entorno separado; verificar FFmpeg, subtítulos, Whisper, revisiones de Nomic/CLAP
y compatibilidad con vectores históricos. Después medir relevancia humana,
latencia y memoria reales. Mantener pendientes P04, P10, P12 y P16 donde corresponda.
\item[E19, P06/P19: modalidades opcionales.] Diseñar disponibilidad por canción e
índices con mapas independientes para admitir instrumentales sin letras ficticias.
Revisar normalización y explicación de pesos antes de implementarlo.
\end{description}
Primero se valida E18 en un entorno reproducible; E17 precede al uso sostenido y
E19 a anunciar ingesta general de instrumentales. Ninguno se considera cerrado
por haber añadido la infraestructura de este bloque.

\section{Corrección de pendientes tras la auditoría de bloques I--IV}
\label{sec:correcciones_pendientes}
La auditoría reprodujo dos debilidades que las 140 pruebas originales no
detectaban: aceptar una traducción equivocada como explícita y aprobar con un
veredicto antiguo sin citas. Se conservan las capturas históricas; las nuevas
pruebas y resultados se guardan en \code{validation/pending_fixes/}.

\subsection{P07--P08/E13: correspondencia de restricciones y abstención}
Antes bastaba que una cita existiera y contuviera una negación: ``sin romance''
podía vincularse al descriptor \emph{violence}. Ahora \code{intent.py} exige
igualdad de frase normalizada o una equivalencia explícitamente revisada. Las
equivalencias son específicas de modalidad y no se aceptan porque el modelo
declare que son correctas. No se eliminan acentos para comparar evidencia.
Una traducción desconocida, ampliación no justificada o inversión de polaridad
provoca error de validación y aclaración mediante el fallback seguro. También
se revalidan restricciones explícitas heredadas de un plan previo.

Se exige que una cita negativa incluya su negación al comienzo y conserve todos
los conceptos de una conjunción. Una cita corta seguida de continuación de lista
no acredita el alcance completo. Se rechazan dobles negaciones, construcciones
``no solo'' y contradicciones entre restricciones explícitas. Los casos
acústicos completamente reconocidos mantienen su gramática determinista.
El prompt pide descriptores mínimos y aclaración cuando no pueda verificarse
una traducción. El control de alcance no pretende resolver toda la gramática.

\textbf{Por qué es mejor.} Se elimina el camino demostrado de etiquetar una
traducción incorrecta como evidencia verificada. La comprobación local no depende
de que otro modelo confirme al primero. \textbf{Coste y alternativa.} Se pierde
cobertura: paráfrasis válidas no incluidas pueden pedir aclaración. Es preferible
a buscar ocultando una restricción; un verificador semántico más amplio sólo debe
añadirse con corpus anotado, tasas de error y abstención medidos. E13 queda
corregido para este contrato conservador; la fidelidad general sigue pendiente
de evaluación humana P04/P10/P12.

\subsection{P10: evidencia uniforme para decisiones del juez}
El cliente ya no registra herramientas de exclusión. Sólo se aplican decisiones
JSON estructuradas con motivo y citas verificables dentro del texto suministrado.
Tanto aprobar como rechazar requieren evidencia no vacía, incluso si la petición
contiene únicamente negativos. Una cita prueba presencia textual, no la
corrección semántica ni por sí sola la ausencia de un tema en toda la canción.
La cobertura incompleta sigue impidiendo aprobar la letra completa.

Un texto antiguo o una llamada de herramienta provoca un único reintento de
formato; si persiste, el juez se abstiene. Si las citas faltan o son inventadas
en ambos intentos, también se abstiene. Los fallos de transporte y otras salidas
malformadas conservan estado de error, sin multiplicar llamadas al proveedor.
El orquestador convierte además una decisión marcada \texttt{legacy} en abstención:
no puede convertirse en aprobación seleccionada ni rechazo semántico.

Las pruebas antiguas de ranking se actualizan para emitir decisiones JSON con
citas; las de compatibilidad comprueban ahora abstención, aislamiento y ausencia
de mutación. Las nuevas regresiones verifican recuperación tras reintento,
negativos sin evidencia, herramientas antiguas, traducción incorrecta, negaciones
compuestas y que una interpretación incierta no llegue a la búsqueda.
El coste posible es más abstenciones y una segunda llamada en respuestas antiguas;
el máximo sigue siendo dos por candidato y el refinamiento continúa desactivado
por defecto. Las citas no sustituyen la calibración humana del juicio.

\subsection{P04/P10/P12: evaluación humana preparada, sin etiquetas inventadas}
\code{evaluation/annotations.py} exporta consultas, metadatos y letras del snapshot
con grados y violaciones vacíos. Para audio exige escucha: no permite inferir
relevancia sonora a partir de la letra. La importación exige declaración de
revisión humana, anotador, consultas idénticas y generación/hashes coincidentes;
valida IDs, grados y cobertura. Esa declaración es atribución, no prueba automática
de que una persona revisó correctamente. Los archivos existentes no se sobrescriben.
\code{evaluation.run --qrels} permite usar la revisión sin modificar el archivo
original; una generación distinta bloquea la evaluación.
Se conserva la distinción entre desconocido y grado cero. No se fabricaron
anotaciones, por lo que relevancia y calibración permanecen abiertas.

\subsection{P16: memoria del proceso y perfil real ejecutable}
El perfil de control ahora registra working set RSS mediante la API de Windows
(o \code{/proc/self/statm} en Linux), muestras de memoria y pico del proceso
durante su vida. El pico de muestras puede perder aumentos breves; el pico de
vida incluye inicialización. Los contadores CUDA, si están disponibles, cubren
asignaciones Torch, no toda la memoria del controlador. No se importa Torch
exclusivamente para medir un arranque bloqueado. No se añadieron dependencias.

\code{evaluation/real_profile.py} ejecuta el runner real con muestreo de RSS, sin
sustituir modelos por controles. Los pesos sólo se buscan localmente por defecto;
la descarga requiere la opción explícita \code{--allow-model-downloads}.
La ejecución realizada quedó bloqueada por \code{GROQ_API_KEY} ausente. Su memoria
es de arranque, no de inferencia. El perfil de control sí mide RSS real de ese
proceso, pero no permite extrapolar uso de memoria o latencia de modelos reales.

\subsection{Estado de cierre y prioridad restante}
Los dos controles de la auditoría dejan de reproducir los fallos. La suite y las
series simuladas se vuelven a capturar con hashes de fuentes actuales, conservando
los informes de bloques anteriores. La documentación TeX se comprueba
estáticamente; no se declara compilación PDF ni revisión visual nueva de la web.
La condición pendiente es configurar el proveedor localmente y realizar la
anotación humana, después ejecutar y comparar modelos reales antes de activar
refinamiento o afirmar calidad/rendimiento. E17--E19 conservan su prioridad y
ubicación del bloque V; esta corrección no realiza ingesta real ni limpieza de datos.

\section{Trabajo autónomo: mantenimiento, instrumentales e integración real}
\label{sec:autonomo}
\textbf{19--20 de octubre de 2026.} Se implementan las actividades que no necesitan
juicios humanos ni una clave del proveedor. Los informes nuevos se conservan en
\code{validation/autonomous/}; los bloques anteriores permanecen como evidencia
histórica. El catálogo del usuario conserva sus 258 canciones: las ingestas de
prueba se realizan en catálogos aislados.

\subsection{E17/P17/P20: cuotas, retención y restauración}
Antes se retenían indefinidamente generaciones, trabajos y medios. Ahora hay
cuotas configurables: catálogo 512 MiB, ingesta 2 GiB y copia 2 GiB por defecto.
La admisión reserva conservadoramente 300 MiB por trabajador activo y por el
siguiente trabajo; la cola mantiene su límite de 32 trabajos. Una publicación
comprueba el uso antes y después de escribir y no cambia el puntero si excede
la cuota. El guard del trabajo comprueba el uso durante el progreso. No se trata
de una cuota impuesta por el sistema operativo: la descarga puede escribir antes
de que un control la detecte; los límites de vídeo y archivo siguen vigentes.
Los pesos de modelos y las copias se administran aparte del catálogo/ingesta.

\code{python -m maintenance cleanup} sólo muestra el plan. La opción
\code{--apply} caduca versiones con más de 30 días, conservando al menos las tres
más recientes y siempre la actual. Conserva medios referenciados por cualquier
versión retenida; trabajos terminales antiguos y archivos huérfanos pueden
retirarse. Los directorios y medios se mueven a cuarentena dentro de su raíz;
el período de recuperación comienza en ese momento. La eliminación permanente
requiere otra operación explícita, \code{purge-trash --apply}, tras la retención.
Los enlaces a versiones caducadas dejan de funcionar: es una contrapartida
del límite de almacenamiento, no una migración transparente de enlaces antiguos.
No se aplicó una eliminación permanente a datos del usuario.

Las rutas se resuelven dentro de raíces administradas y se rechazan enlaces
simbólicos y junctions. Si una generación retenida está corrupta se aborta el
plan, en lugar de suponer que sus medios son prescindibles. El mantenimiento
requiere una cola sin trabajos activos y toma bloqueos en el mismo orden que
publicación. Se caduca el historial antes de mover archivos: una caída puede
dejar archivos sobrantes; no cambia el puntero actual para apuntar a archivos
eliminados. Las dos bases SQLite no constituyen una transacción distribuida.

\code{backup --path} crea un ZIP con catálogo publicado, medios referenciados y
ambas bases SQLite serializadas; excluye claves, configuración privada, pesos,
trabajos temporales y cuarentena. Verifica SHA-256 de cada entrada antes de
publicar el ZIP. \code{restore --path --destination} exige una carpeta nueva,
valida inventario, hashes, integridad SQLite, índices y presencia de medios;
no sobrescribe ni cambia automáticamente el servicio en ejecución. Hay una copia
real del catálogo actual y una restauración verificada en una carpeta aislada.
Para una recuperación se configuran ambas rutas restauradas y se reinicia el
servicio. Los hashes detectan corrupción; no autentican un ZIP de un tercero.

\subsection{E19/P06/P19: instrumentales sin letras ficticias}
El mapa v1 imponía una fila por posición en ambos índices. El mapa v2 conserva
un mapa completo de catálogo y mapas independientes para letras y audio. Se
validan posiciones, IDs únicos, disponibilidad del texto y cobertura acústica.
Las generaciones v1 siguen siendo legibles, sin reconstruir sus vectores.
Una canción sin letra no tiene posición ni vector en el índice lírico; se
mantienen audio y metadatos. Los workers reciben el mapa de su modalidad dentro
del snapshot fijado, evitando atribuir a una canción el vector de otra.

La API admite \code{instrumental:true}; la interfaz ofrece una casilla con esa
decisión. Se omiten subtítulos, Whisper y Nomic para ese trabajo y se registra
que la condición instrumental fue declarada, no verificada acústicamente. Si no
se declara instrumental, una transcripción vacía continúa produciendo error.
La búsqueda acústica puede recuperar el registro; una búsqueda lírica no lo
recupera desde el índice de letras. Un candidato sin texto no recibe aprobación
lírica por el mero hecho de coincidir en audio. La tarjeta explica la ausencia
de letra y no presenta el sonido como auditado.

\subsection{E18/P01/P19: entorno y modelos comprobados realmente}
Se instala \code{.venv-ingestion} separado del entorno histórico. Se resuelven
las versiones completas en \code{requirements/ingestion.lock}; el cierre de
dependencias de modelos se actualiza en \code{requirements/models.lock}.
FFmpeg se aprovisiona desde \code{imageio-ffmpeg} a una ruta local explícita y
no se cambia el PATH global. Whisper recibe muestras decodificadas a 16 kHz;
CLAP usa muestras a 48 kHz. Los caches de pesos quedan en \code{.tools/}.
\code{pip check}, las importaciones y las regresiones se comprueban en el entorno.

La primera inferencia real descubrió una incompatibilidad de Nomic con
Transformers 5: faltaba \code{get_extended_attention_mask}. La solución limita
Transformers a la serie 4 compatible, con 4.57.6 verificada. No se parchea código
descargado ni se oculta la excepción. Se fijan revisiones de Nomic, su código
remoto y CLAP en \code{requirements/model-revisions.json}. Las publicaciones
registran esas revisiones locales; no se atribuyen retroactivamente al corpus
heredado, cuya procedencia exacta sigue sin demostrarse.

La ingesta técnica usa un vídeo de prueba de yt-dlp: descarga real, conversión
MP3, tres ventanas CLAP y publicación aislada con 259 canciones, 258 vectores
líricos y 259 acústicos. Nomic genera un vector normalizado de 256 dimensiones
y Whisper base procesa dos segundos de silencio generado en CPU. El silencio
comprueba ejecución, no exactitud de transcripción de canciones. Los workers
reales ejecutan consultas e índices con muestras iniciales y tres muestras
calientes; éstas pueden usar caché y no equivalen a nuevas inferencias.
Se registra RSS del proceso. No se mide rendimiento del router/juez ni calidad
musical a partir de estas comprobaciones.
El vídeo antiguo estaba retirado; la sustitución se documenta en el
\href{https://github.com/yt-dlp/yt-dlp/issues/12263}{repositorio de yt-dlp}.

\subsection{Interfaz y documento}
La prueba de navegador usa cola, publicación y FFmpeg reales con IA de control,
en un catálogo aislado y con un aviso visible. Comprueba progreso, búsqueda
simultánea, actualización de conteos, recuperación del último trabajo tras
recargar y tarjeta instrumental con reproducción. Se guardan capturas y se
corrige la presentación de la casilla y la conservación de años sin decimales.
La compilación PDF utiliza MiKTeX; se corrigen saltos de línea de rutas técnicas,
se resuelven referencias y se revisan páginas renderizadas. El informe de
compilación distingue esto de las comprobaciones estáticas anteriores.

\subsection{Lo que todavía necesita participación del usuario}
Configurar \code{GROQ_API_KEY} localmente para evaluar router, juez y refinador
reales; no compartir la clave en la conversación. Completar y atribuir las
anotaciones humanas, escuchar las muestras acústicas y revisar las transcripciones
de música antes de declarar relevancia o fidelidad. Después se compara el
refinamiento y se decide si su beneficio compensa coste y latencia. El acceso
a Groq y los juicios humanos no se sustituyen por inferencias de control.
La política de retención es configurable y conservadora; puede ajustarse al uso
real sin que eso impida entregar las herramientas y pruebas de mantenimiento.

\subsection{Compatibilidad del proveedor durante la prueba local}
El 20 de octubre de 2026 se verificó que la credencial permite consultar Groq, pero el modelo Scout configurado devuelve \texttt{model\_not\_found}. Se seleccionó en la configuración local \texttt{openai/gpt-oss-120b}, incluido en los modelos accesibles de la cuenta. También se explicitaron los cuatro nombres de campos obligatorios del planificador y la obligación de conservar las listas vacías: la respuesta anterior empleaba un nombre alternativo incompatible. La consulta ``triste bajo la lluvia produjo un plan válido con estado \texttt{ready}. Esta prueba confirma conexión y contrato, pero no demuestra calidad musical ni valida el flujo completo de búsqueda. Queda pendiente comparar calidad, latencia y coste del modelo con consultas revisadas por personas.


\section{Registro crítico de cambios realizados}
\label{sec:evolucion}
Este registro documenta la sustitución del estado inicial A/B/C por el flujo
actual. Los ocho números corresponden a los dos bloques de implementación
acordados. Los puntos E01--E12 conservan el diagnóstico adicional. Su ejecución y la de
los antiguos puntos 9--14 se organiza ahora mediante P01--P20 en la
sección~\ref{sec:plan}; no son tres listas de trabajo independientes.

\subsection{Cómo se juzga una mejora}
Se distinguen cuatro tipos de conclusión:
\begin{description}
\item[Corrección comprobada:] una prueba reproduce una conducta incorrecta y
comprueba la propiedad esperada en el código actual.
\item[Optimización comprobada por conteo:] se comprueba que se realizan menos
cargas o llamadas equivalentes. Esto no acredita un porcentaje de reducción de
latencia, memoria total o coste real.
\item[Capacidad añadida:] existe una ruta de ejecución antes ausente. Su utilidad
para el usuario necesita evaluación funcional y de relevancia.
\item[Decisión con contrapartidas:] mejora una propiedad y puede empeorar otra.
Debe registrarse la alternativa y su condición de aceptación.
\end{description}
La evidencia disponible comprende 32 pruebas aprobadas al finalizar el bloque 8,
con LangGraph real y proveedores/modelos simulados, y análisis sintáctico de 17
archivos activos y de pruebas. Se probaron concurrencia del juez e inicialización
de recursos; no se hizo una prueba de carga HTTP con modelos reales.

\subsection{Cambio previo: nombre e identidad del proyecto}
Se sustituyeron las referencias textuales a la marca anterior por SPELLFIND en
12 archivos, incluidos interfaz, prompts, documentación y notebook; también se
renombró el archivo del grafo a \code{grafo_SPELLFIND.png}. Se conservaron los
identificadores de datasets externos y los títulos bibliográficos históricos.

La ventaja es la consistencia de nombres y referencias locales. No cambia el
algoritmo ni demuestra mayor calidad de recuperación. El renombrado del PNG no
implica que se haya redibujado o actualizado el contenido de la imagen.

\subsection{Punto 1: exclusiones persistentes}
\label{cambio:1}
\textbf{Antes.} A repetía la búsqueda final con un conjunto vacío de exclusiones,
de modo que una canción rechazada podía reaparecer. C añadía también las canciones
aprobadas a sus exclusiones.

\textbf{Ahora.} El grafo conserva un pool por solicitud y filtra el ranking al
seleccionar candidatos y al generar la respuesta. Solo los rechazos se agregan
al pool; el candidato aprobado queda primero. Si faltan candidatos, no se
reincorporan rechazados para completar cinco resultados. Se conserva el motivo
de cada exclusión. \code{ExclusionPool.exclude} evita duplicar su registro cuando
el juez y el orquestador registran el mismo rechazo.

\textbf{Archivos.} \code{judge.py}, \code{workflow.py} y \code{app.py}.
\textbf{Evidencia.} Las pruebas de rechazos integrados y de agotamiento verifican
que los IDs excluidos no vuelven, incluso después de refinar consultas.

\textbf{Balance.} Corrección de coherencia comprobada. La política conserva
también rechazos posiblemente equivocados del LLM; aumentar su precisión es
distinto de hacer cumplir sus decisiones. Un candidato rechazado bajo una
consulta anterior no se reconsidera después de refinar. E01 y E06 revisan esa
política y la calidad de la evidencia.

\subsection{Punto 2: veredictos explícitos y errores visibles}
\label{cambio:2}
\textbf{Antes.} A calculaba una variable de aprobación, pero devolvía verdadero
ante cualquier respuesta sin herramientas. Los errores del proveedor también
se convertían en aprobación.

\textbf{Ahora.} El juez devuelve \code{approved}, \code{rejected} o \code{error}.
Solo acepta aprobación textual explícita. Un rechazo textual reconocido o una
llamada válida a \code{exclude_song} rechaza al candidato. Respuestas vacías,
ambiguas, argumentos inválidos y fallos del proveedor producen error sin alterar
el pool. Todas las llamadas de herramienta se validan antes de excluir, incluido
que el ID coincida con el candidato actual.

Se eliminó la segunda llamada al LLM que solo confirmaba el rechazo: la exclusión
se resuelve después de validar la primera respuesta. La interfaz muestra una
advertencia cuando no se completa la validación. La API conserva el booleano
\code{approved} por compatibilidad y añade estado explícito y estado por resultado.

\textbf{Archivos.} \code{judge.py}, \code{app.py} y \code{templates/index.html}.
\textbf{Evidencia.} Pruebas de texto ambiguo, rechazo textual, herramientas
malformadas, ID ajeno y fallos de inicialización/proveedor.

\textbf{Balance.} Se corrigen falsas aprobaciones por errores de control. Esto no
demuestra que el LLM juzgue mejor una letra. El parseo estricto puede convertir
una respuesta semánticamente válida pero distinta del formato en error; no hay
reintento limitado ni validación humana del juez. Pendiente E06.

\subsection{Punto 3: contrato único de planes y fallback}
\label{cambio:3}
\textbf{Antes.} El prompt pedía listas, mientras que el fallback devolvía un
diccionario lírico incompatible con el worker.

\textbf{Ahora.} \code{search_contract.py} valida las cuatro listas. Adapta cadenas
individuales y el formato lírico antiguo con \code{text} y \code{weight}, elimina
espacios, textos vacíos y duplicados exactos, y rechaza elementos no textuales,
campos ausentes o planes sin consultas positivas. Los nulos se convierten en
listas vacías. El peso antiguo se valida pero no determina la fusión: los pesos
se siguen infiriendo en la aplicación. Router, fallback y refinador comparten
este contrato. El grafo puede usarlo incluso si falla la importación del router.

\textbf{Archivos.} \code{search_contract.py}, \code{router.py},
\code{workflow.py} y \code{refiner.py}.
\textbf{Evidencia.} Pruebas de esquemas antiguos, nulos, listas mixtas, duplicados,
falta de positivas y llegada del fallback a los workers y al juez.

\textbf{Balance.} Se comprueba compatibilidad de tipos, no calidad semántica.
El fallback copia la consulta original en ambas modalidades, sin traducir al
inglés ni preservar negativos detectados por un plan fallido. Puede empeorar
la recuperación y activar una modalidad innecesaria. Pendiente E02.

\subsection{Punto 4: aislamiento por solicitud}
\label{cambio:4}
\textbf{Antes.} A guardaba el pool activo en \code{_active_pool}; dos solicitudes
intercaladas podían modificar el pool de la otra.

\textbf{Ahora.} Se elimina esa variable. El pool se pasa explícitamente al juez;
la herramienta describe el rechazo y la aplicación lo aplica al contexto local
tras validarlo. Cada búsqueda tiene \code{request_id}, logs y copias al consultar
el conjunto de exclusiones. Una búsqueda nueva inicia sin rechazos previos.

\textbf{Archivos.} \code{judge.py}, \code{workflow.py} y \code{app.py}.
\textbf{Evidencia.} Una prueba con barrera intercala dos evaluaciones y verifica
que cada pool contiene solo su propio ID; otra ejecuta búsquedas sucesivas.

\textbf{Balance.} Elimina el cruce de exclusiones reproducido. No certifica que
todos los clientes externos o modelos soporten cualquier nivel de concurrencia.
Además, el pool se muta desde el juez y desde el grafo; la idempotencia evita
duplicados, pero conviene asignar un único propietario del estado. Pendiente E08.

\subsection{Punto 5: consolidación en un único grafo}
\label{cambio:5}
\textbf{Antes.} La web híbrida y el grafo semántico tenían reglas diferentes.
B era un borrador con incompatibilidades y C no recuperaba audio.

\textbf{Ahora.} \code{workflow.py} contiene el flujo común con estado tipado y
dependencias inyectadas: planificar, recuperar, evaluar y refinar. Web y consola
lo invocan. \code{cli.py} acepta una consulta por argumento. Los nombres antiguos
\code{tarea_3.py} y \code{tarea_3_correcion.py} redirigen a esa CLI y requieren
también un argumento; ya no ejecutan las consultas de ejemplo embebidas.
Los originales se conservaron en \code{experiments/}, marcados como históricos.

\textbf{Ventaja.} Hay una sola definición de exclusión, límite, refinamiento y
error para las entradas soportadas. La inyección permite ejecutar el grafo real
con modelos simulados en pruebas. Se verifican por inspección las entradas
compartidas; no se ejecutó una sesión CLI completa con modelos reales.

\textbf{Contrapartida.} LangGraph pasa a ser dependencia del flujo web; el grafo
se compila de nuevo por solicitud. Frente a C, el flujo nuevo no conserva
\code{MemorySaver}, streaming de actualizaciones por CLI ni recuperación de un
checkpoint tras una excepción. Por ello es una consolidación de lógica, no un
superconjunto de todas las capacidades previas. Pendiente E04.

\subsection{Punto 6: refinamiento selectivo por modalidad}
\label{cambio:6}
\textbf{Antes.} A descartaba candidatos con consultas fijas. C reescribía solo
letras y acumulaba vectores, con riesgo de repetir o arrastrar consultas fallidas.

\textbf{Ahora.} Después de un rechazo, \code{refiner.py} recibe consulta original,
plan, candidato y motivo del juez. Propone un plan nuevo y las modalidades que
necesitan cambiar. La validación impide cambiar las otras listas, eliminar
negativos existentes, vaciar una modalidad activa o activar una inicialmente
ausente. La normalización elimina duplicados exactos. Se admiten cambios
líricos, acústicos o de ambos; un plan sin cambios también es válido.

El prompt pide conservar imágenes literales, reemplazar positivos en vez de
acumularlos e impedir restricciones inventadas. Estas son instrucciones al
modelo: el validador no puede certificar por sí solo que una frase nueva conserve
la intención. El plan inicial, el final y cada refinamiento quedan registrados.

\textbf{Archivos.} \code{refiner.py}, \code{workflow.py} y \code{app.py}.
\textbf{Evidencia.} Pruebas de cambios por modalidad, planes sin cambios,
conservación de negativas, cambios indebidos y fallo del refinador.

\textbf{Balance.} Capacidad añadida, con mejora de relevancia aún no demostrada.
Un rechazo lírico no demuestra un fallo acústico: el refinador declara que no
hay evidencia de audio. Puede reescribir descriptores a partir de la petición,
pero no juzgar cómo suena el candidato. También añade una llamada LLM después
de cada rechazo refinable, hasta dos con tres intentos. Si falla, hoy se detiene
la búsqueda aunque queden candidatos del plan anterior. Pendientes E01, E03,
E05 y E07.

\subsection{Punto 7: terminación con significado explícito}
\label{cambio:7}
\textbf{Antes.} C enviaba a la ruta denominada aceptación al alcanzar tres
veredictos incluso si todos eran insuficientes. A podía devolver resultados
sin una aprobación clara al terminar el ciclo.

\textbf{Ahora.} \code{approved} requiere un veredicto explícito;
\code{limit_reached} significa presupuesto agotado con candidatos restantes;
\code{no_candidates} significa que no quedan candidatos recuperados sin excluir;
\code{error} distingue fallos y añade la etapa. No se refina después del último
intento ni se excluye una canción aprobada. Si el ranking queda vacío, el juez
no se invoca. Una canción sin letra consume un intento y se excluye sin LLM.

\textbf{Archivos.} \code{workflow.py}, \code{app.py}, interfaz y CLI.
\textbf{Evidencia.} Pruebas de límite, agotamiento, errores y ausencia de letras.

\textbf{Balance.} Corrección de control comprobada. El presupuesto cuenta
selecciones de candidato, no solo llamadas al juez. La ausencia de letra no
debería equivaler siempre a fallo de una consulta acústica o instrumental.
\code{no_candidates} describe el ranking explorado, no demuestra que todo el
catálogo carezca de canciones válidas. Pendientes E01 y E10.

\subsection{Punto 8: reutilización de rankings, embeddings y recursos}
\label{cambio:8}
\textbf{Antes.} A repetía ambas búsquedas hasta tres rondas y una búsqueda final,
aunque las consultas fueran iguales. Servidor y workers cargaban copias de
índices, mapa y CSV; los negativos se recodificaban para cada positiva.

\textbf{Ahora.} Cada solicitud cachea resultados por modalidad, positivas,
negativas y top-K. Las exclusiones se aplican después, por lo que no invalidan
esa caché. Cambiar solo letras reutiliza audio y viceversa. El resultado final
usa el último ranking disponible sin otra recuperación.

\code{resources.py} comparte catálogo, mapa e índices por proceso, con
inicialización protegida. El lookup usa un diccionario por ID y devuelve una
copia de los metadatos. El catálogo valida IDs duplicados y normaliza valores
ausentes. Los workers conservan sus modelos bajo bloqueo y no leen el CSV.
Cada uno tiene una LRU de 512 embeddings, devuelve copias y protege la
codificación para impedir mutaciones compartidas y cargas concurrentes dobles.

\textbf{Evidencia.} Con tres evaluaciones y consultas sin cambios, las pruebas
comprueban una llamada a cada worker frente a las hasta cuatro de A. Con un
cambio de una modalidad, comprueban dos llamadas a esa rama y una a la otra.
Otras pruebas verifican carga única de recursos/modelo bajo concurrencia y
copias independientes de embeddings.

\textbf{Balance.} Menos trabajo redundante demostrado por conteo; no existe aún
un benchmark con modelos reales. Las caches consumen memoria y no sobreviven
al proceso. Se comparte por proceso, no entre procesos de servidor. Se mantienen
tanto DataFrame como diccionario de metadatos, por lo que no se ha eliminado
toda duplicación de representación. Los bloqueos serializan codificaciones
por modalidad y el catálogo no se actualiza en caliente. Pendientes E07 y E09.

\subsection{Cambios auxiliares y compatibilidad}
\begin{itemize}
\item La API añade \code{termination_reason}, \code{error_stage},
\code{initial_plan}, \code{refinement_log}, \code{workflow_trace} y
\code{validation_scope: lyrics}. Se mantienen los campos de resultados previos.
\item Se devuelven \code{n_lyrics} y \code{n_audio}, que la interfaz ya intentaba
mostrar. Son cantidades del último ranking por rama, no contadores acumulados
de operaciones o canciones únicas de todo el proceso.
\item El servicio valida consulta no vacía y k positivo. El endpoint aún llama
a métodos del cuerpo antes de validarlo completamente: no está resuelto el
manejo HTTP de cualquier tipo de JSON. Pendiente E11.
\item Las rutas de lectura de catálogo e índices se resuelven desde el módulo.
El router y el indexador todavía mantienen rutas relativas al directorio de
trabajo; no se declara independencia completa de ese directorio.
\item Los workers ahora aceptan conceptos de cinco caracteres o menos y filtran
solo vacíos. Evita perder términos como piano, pero su impacto en relevancia no
ha sido medido.
\item El flujo actual propaga fallos de recuperación y descarta un ranking
antiguo si falla una búsqueda tras refinar. Evita mostrar scores asociados al
plan incorrecto. Frente a A se pierde la devolución parcial si una sola rama
falla; esa política debe revisarse, pendiente E05.
\item Se añadieron README, documentación, pruebas y exclusión de temporales en
Git. No se reescribieron los datos del catálogo ni se midió una fórmula nueva
de fusión: el score porcentual sigue pendiente de corrección.
\end{itemize}

\section{Puntos extra por realizar o corregir}
\label{sec:extras}
Esta lista conserva el diagnóstico previo al bloque I. E10 tiene cubierta su validación numérica por P03 y mantiene presupuestos pendientes en P11; E11 tiene entorno y HTTP verificados en P01--P02, con validación integral pendiente en P20; E07 tiene protocolo y simulación en P04, con calidad real pendiente. Los demás puntos continúan pendientes. El orden de esta lista
es histórico; la prioridad de ejecución vigente es la del plan P01--P20.
Se registran alternativas
con una condición verificable, no como mejoras ya implementadas. Prioridad alta
significa riesgo de respuesta incorrecta o pérdida de disponibilidad; media
indica decisiones de rendimiento, operación o mantenibilidad por medir.

\subsection{E01. Separar validación lírica, acústica y de consultas instrumentales}
\textbf{Prioridad alta.} El juez sigue leyendo letras y el grafo excluye canciones
sin letra. Una consulta solo acústica puede necesitar un resultado instrumental
que hoy se descartaría; una aprobación no valida audio. Además, un rechazo
lírico no justifica por sí solo refinar CLAP.

\textbf{Acción.} Definir requisitos y evidencia por modalidad; emitir
\code{not_evaluated} cuando falte evidencia acústica, permitir rutas solo audio,
y seleccionar refinamientos a partir de feedback fundamentado. Evaluar la
reconsideración de rechazos dependientes de una consulta refinada manteniendo
fijas las restricciones explícitas del usuario.
\textbf{Aceptación.} Consultas instrumentales no se rechazan por ausencia de
letra y ninguna salida afirma validación acústica sin evidencia.

\subsection{E02. Fallback fiel al idioma, modalidades y restricciones}
\textbf{Prioridad alta.} El respaldo usa el texto original en ambas ramas; su
formato es correcto, pero no garantiza traducción, separación modal ni negativos.
\textbf{Acción.} Diseñar un fallback degradado que conserve el último plan válido
o restricciones verificables e informe qué interpretación no pudo hacerse.
No inventar una extracción fiable de negativos mediante reglas sin evaluarla.
\textbf{Aceptación.} Consultas con negaciones e idiomas distintos del inglés no
pierden silenciosamente sus restricciones cuando falla el router.

\subsection{E03. Proveniencia de negativos y activación de modalidades}
\textbf{Prioridad alta.} Conservar todos los negativos impide quitar uno
inventado por el LLM. Impedir activar modalidades ausentes bloquea la reparación
de un error inicial del router.
\textbf{Acción.} Separar restricciones explícitas del usuario de hipótesis del
modelo; bloquear las primeras y permitir revisar las segundas mediante evidencia.
\textbf{Aceptación.} Se puede reparar una clasificación modal incorrecta sin
eliminar restricciones explícitas ni añadir requisitos no pedidos.

\subsection{E04. Coste del grafo y funciones perdidas frente a C}
\textbf{Prioridad media.} El grafo se compila por solicitud y ya no hay checkpoints
ni streaming equivalentes a C. LangGraph no hace que la recuperación sea mejor
por sí mismo.
\textbf{Acción.} Medir el coste de compilación y valorar un grafo compilado
reutilizable con estado por solicitud. Si se necesita reanudar o transmitir
progreso, recuperar esas capacidades con contratos serializables. Comparar con
una máquina de estados simple si el grafo no aporta ventajas medibles.
\textbf{Aceptación.} Igualdad funcional de entradas, aislamiento concurrente y
medidas de coste; declarar explícitamente si se ofrece reanudación o streaming.

\subsection{E05. Recuperación parcial y fallos del refinador}
\textbf{Actualización: resuelto técnicamente por P09 en la
sección~\ref{sec:bloque3}. El diagnóstico siguiente corresponde al estado previo.}
\textbf{Prioridad alta.} Ahora un fallo de una rama detiene toda la recuperación;
un fallo del refinador impide evaluar el siguiente candidato ya disponible.
Son decisiones conservadoras, no una mejora automática de disponibilidad.
\textbf{Acción.} Diseñar un estado degradado por modalidad y una salida que
conserve evidencia del plan usado. Evaluar continuar con los rankings previos
cuando falla el refinador y reajustar pesos solo con una política validada.
\textbf{Aceptación.} Un fallo de CLAP no borra resultados líricos útiles sin
explicación y un candidato no se marca aprobado por omitir la rama fallida.

\subsection{E06. Calidad del juez, formato y extensión de la evidencia}
\textbf{Actualización: P10 implementa formato, citas, cobertura y reintento;
calibración humana pendiente. Véase~\ref{sec:bloque3} para el estado actual.}
\textbf{Prioridad alta.} El formato estricto evita aprobaciones accidentales,
pero no calibra el juicio. Solo se inspeccionan 3000 caracteres, se juzga el
top-1 y no hay reintento de formato. Los demás resultados quedan sin evaluar.
\textbf{Acción.} Comparar salida estructurada, reintento limitado y políticas
de evidencia por fragmentos con referencia humana. Mantener fallos como errores
y diferenciar abstención de rechazo. Evaluar rechazo persistente y revisión
después de refinar sobre ejemplos etiquetados.
\textbf{Aceptación.} Medir falsos rechazos y aprobaciones, cubrir pistas situadas
al final de la letra y no presentar todo el top-K como validado por evaluar uno.

\subsection{E07. Comprobar si refinar y cachear mejora el resultado global}
\textbf{Actualización: P12 tiene comparación simulada y política experimental;
la utilidad musical y las medidas reales siguen pendientes. Véase~\ref{sec:bloque3}.}
\textbf{Prioridad alta para afirmar superioridad.} Las pruebas verifican ramas y
conteos; no prueban más relevancia ni menor latencia de extremo a extremo.
El refinador añade llamadas y puede desviar la intención original.
\textbf{Acción.} Comparar el mismo baseline corregido con y sin refinamiento,
y medir latencia fría/caliente, llamadas, coste, memoria, relevancia y negativos.
Usar el mismo corpus, consultas y criterios; repetir y reportar dispersión.
\textbf{Aceptación.} Publicar una comparación reproducible y habilitar el
refinamiento solo donde su beneficio justifique coste y errores adicionales.

\subsection{E08. Un único propietario de exclusiones y estado serializable}
\textbf{Prioridad media.} Juez y grafo mutan el pool; la idempotencia evita
duplicados, pero el estado contiene objetos y caches que no se han validado para
checkpoints. \code{TypedDict} no impone validación en tiempo de ejecución.
\textbf{Acción.} Hacer que el juez devuelva decisiones y el orquestador sea el
único que aplique exclusiones. Separar estado serializable y recursos efímeros,
con invariantes de aprobación, rechazo y versión del plan.
\textbf{Aceptación.} Una exclusión tiene un solo punto de escritura; si se añade
persistencia, guardar y restaurar conserva las decisiones sin compartir pools.

\subsection{E09. Invalidación de caches, memoria y concurrencia real}
\textbf{Actualización: P16 instrumenta contención y cachés y prueba concurrencia
local con índices reales/modelos de control; invalidación y memoria real siguen
pendientes de P16/P17. Véase~\ref{sec:bloque4}.}
\textbf{Prioridad media.} Reiniciar es hoy la política de actualización. Índices,
modelos, DataFrame y diccionario consumen memoria por proceso; devolver un
DataFrame o mapa compartido permite que otro consumidor lo modifique.
\textbf{Acción.} Medir RSS, GPU y contención; valorar snapshots versionados,
lecturas inmutables, publicación atómica del catálogo, límites de caché medidos
y batch de embeddings. No eliminar bloqueos sin comprobar el modelo concreto.
\textbf{Aceptación.} Resultados coherentes con una versión del índice, consumo
acotado y carga concurrente representativa; nunca mezclar mapa e índice viejos
y nuevos.

\subsection{E10. Presupuesto de candidatos y validación numérica}
\textbf{Actualización: validación numérica cubierta por P03 y presupuestos
separados por P11; véanse~\ref{sec:bloque1} y~\ref{sec:bloque3}.}
\textbf{Prioridad alta.} El límite cuenta también candidatos sin letra y solo
se exploran rankings de tamaño fijo. Faltan controles numéricos completos de
normas, valores finitos y compatibilidad índice--embedding. La resta negativa
puede producir una dirección nula o poco informativa.
\textbf{Acción.} Separar intentos de evaluación, descartes de metadatos y
ampliación de recuperación; validar vectores y catálogo al cargar.
\textbf{Aceptación.} Un fallo de datos no consume todo el presupuesto útil;
una consulta degenerada no entra a FAISS como una búsqueda aparentemente válida.

\subsection{E11. Entorno reproducible y contrato HTTP completo}
\textbf{Prioridad alta.} No se ha comprobado una instalación limpia de todas las
dependencias fijadas. El entorno temporal de pruebas recibió instalaciones
superpuestas; la última salida de metadatos informó LangGraph 1.2.16 y
LangChain Core 1.4.3 aunque también se intentó instalar LangGraph 1.2.4. Esa
salida no certifica un lock limpio ni la matriz completa de compatibilidad.
\textbf{Acción.} Crear un entorno aislado reproducible y separar dependencias
directas y de desarrollo. Verificar imports reales, versiones y resolución;
probar HTTP con Flask real, JSON inválido, estados 400/500/503 e inicialización
WSGI. Completar rutas, configuración, guía de instalación y modo de despliegue.
\textbf{Aceptación.} Instalar y ejecutar desde cero con versiones inequívocas,
sin depender de la carpeta temporal de esta sesión.

\subsection{E12. Completar los puntos originales 9--14}
\textbf{Actualización: P13--P15 cubren puntuación, pesos, filtros y evidencia
visible, con límites de calidad explícitos. P16 registra un perfil local; carga
real e ingesta P19 siguen pendientes. Véase~\ref{sec:bloque4}.}
\textbf{Prioridad según alcance.} Siguen pendientes puntuación interpretable y
pesos ajustables, filtros exactos, evidencia por resultado, operación completa,
evaluación reproducible e ingesta desde YouTube. El grafo unificado no los da
por implementados. E06, E07 y E11 detallan trabajos relacionados, sin duplicar
la numeración acordada.
\textbf{Aceptación.} Cada bloque futuro identifica los puntos cubiertos y sus
pruebas, y actualiza este registro sin atribuir mejoras todavía no medidas.

\subsection{E13. Fidelidad semántica de citas y cobertura del respaldo}
\textbf{Prioridad alta. Estado pendiente tras el bloque II.} La procedencia
verifica una cita literal y su negación, pero no demuestra que el texto inglés
propuesto por un modelo tenga exactamente el mismo significado. Los controles
léxicos tampoco cubren todas las lenguas, negaciones o contrastes. Una hipótesis
mal traducida podría quedar inmovilizada si se etiqueta como explícita.
\textbf{Acción.} Evaluar equivalencia de restricciones con ejemplos anotados,
separar evidencia textual de confirmación semántica y medir cobertura y coste
de reformulación. Considerar entrada estructurada de restricciones cuando aporte
una decisión útil al usuario, sin certificar una traducción solo por citarla.
\textbf{Aceptación.} Medir pérdidas/adiciones de intención y tasa de abstención;
ampliar reglas o modelos solo con evidencia. Se integra en P10/P12, y en P15
si se decide exponer corrección de restricciones al usuario.

\subsection{E14. Versionado por cambios de orden sin cambio de intención}
\textbf{Prioridad alta. Diagnóstico histórico; corregido el 18 de octubre de 2026.}
La huella del plan depende del orden de sus listas. Un refinador que solo
invierte dos consultas crea nuevas versiones y libera exclusiones previas.
Con dos candidatos y rechazo constante se evaluó al primero tres veces:
IDs $[1,1,1]$, versión final 3 y límite agotado, sin evaluar el segundo.
\textbf{Acción previa al bloque III.} Canonizar la identidad de consultas y
restricciones equivalentes, separando cambios de presentación de cambios
efectivos. Conservar evidencia y restricciones explícitas.
\textbf{Aceptación.} Permutar elementos equivalentes no libera rechazos;
cambios efectivos justificados sí permiten reconsiderar candidatos.
Asignado al bloque II/P07, con regresiones en P11/P12. La corrección y prueba
de no repetición constan en la sección~\ref{sec:bloque3}.

\subsection{E15. Serialización sin validar en la ruta de aclaración}
\textbf{Prioridad media. Diagnóstico histórico; corregido el 18 de octubre de 2026.}
El retorno temprano de planificación que solicita aclaración conserva metadatos
sin comprobar su serialización finita. Una entrada controlada con un campo
numérico NaN alcanzó el estado final de aclaración; JSON estricto lo rechaza.
Esto contradice la garantía general de estado JSON de P05. No se observó una
respuesta real del proveedor con ese dato.
\textbf{Acción previa al bloque III.} Aplicar el mismo contrato a todas las
salidas o definir un esquema explícito de metadatos permitidos.
\textbf{Aceptación.} Probar planes listos y aclaraciones con números no finitos
y objetos no serializables; todos los estados aceptados admiten JSON estricto.
La evidencia de E14/E15 y hashes de fuentes están en
\texttt{validation/revision\_block2/report.json}; son controles locales con
LangGraph real y colaboradores simulados, sin medir relevancia musical.
Asignado al bloque II/P05/P08; ambas rutas ahora validan metadatos JSON finitos.
Las pruebas actuales constan en la sección~\ref{sec:bloque3}.

\subsection{E16. Registro de consola incompatible con CP1252}
\textbf{Prioridad media. Diagnóstico histórico; corregido el 18 de octubre de 2026.}
En esta consola Windows, la suite de 72 pruebas terminó con dos errores
\texttt{UnicodeEncodeError} al imprimir símbolos Unicode desde el router y
el worker lírico durante su inicialización. Ejecutar con
\texttt{python -X utf8 -m unittest discover -s tests} dio 72 pruebas aprobadas.
No se descargaron modelos: las inicializaciones empleaban colaboradores
simulados. El worker CLAP contiene un mensaje similar, identificado por lectura
pero no reproducido en esta ejecución.
\textbf{Acción previa al bloque III.} Usar registros compatibles con el destino
y documentar UTF-8 como alternativa operativa. Un mensaje no debe interrumpir
la inicialización. \textbf{Aceptación.} Verificar salidas CP1252 y UTF-8 sin
fallos causados por el registro. Complementa P01/P02; limita la afirmación de
portabilidad de su verificación sin invalidar los demás controles.
Asignado al bloque I/P01/P02. La continuación introduce salida compatible y
verifica las 96 pruebas actuales en ambas codificaciones; véase~\ref{sec:bloque3}.

\section{Mantenimiento de la documentación y trazabilidad}
\label{sec:mantenimiento}
Cada entrega futura debe actualizar el README operativo y este registro antes de
declararse terminada. La entrada de cada cambio debe incluir fecha, requisito,
archivos, comportamiento anterior y actual, motivo, pruebas ejecutadas,
limitaciones, compatibilidad y tareas adicionales. Los puntos extra deben
mantener su ID y registrar cuándo pasan de pendientes a implementados o
descartados, con evidencia.

La descripción de A/B/C permanece como comparación histórica. Las capacidades
actuales se identifican como flujo consolidado; no se deben corregir los
originales archivados para que parezca que ya tenían las mejoras. Los resultados
experimentales de las presentaciones conservan su fecha y contexto.

Una afirmación de mejor rendimiento requiere medidas; una de mejor calidad
requiere evaluaciones comparables. Las pruebas con modelos simulados acreditan
contratos y control, pero no sustituyen pruebas con Groq, Nomic y CLAP reales.
% Estado adicional conservando el diagnóstico histórico.
\paragraph{18 de octubre de 2026: bloque V.}
Se incorpora publicación transaccional por generaciones, contexto de snapshot
por consulta, cola persistente con propiedad temporal y adaptador opcional de
ingesta. Se evita mezclar recursos o publicar trabajos vencidos; la búsqueda sigue
operativa durante la ingesta. La decisión de no recuperar checkpoints de búsqueda
y los costes de copiar índices se documentan en la sección~\ref{sec:bloque5}.
La calidad real y la compatibilidad de vectores nuevos no están demostradas.
E17--E19 conservan retención, integración real e instrumentales como pendientes.



\appendix
\section{Trazabilidad}
\small
\begin{longtable}{p{0.33\textwidth}p{0.57\textwidth}}
\toprule Concepto & Evidencia\\\midrule\endhead
YouTube, Whisper, db.json e ingesta & Enunciado del usuario.\\
Corpus histórico y limpieza & PDF de mayo, pp. 3--4.\\
Multi-query y RRF & PDF de mayo, pp. 5--6; workers locales.\\
Evaluación 0--3 & PDF de mayo, pp. 8--9.\\
Tres herramientas y roles & PDF de junio, pp. 4--5.\\
Comparación V1/V2 & PDF de junio, p. 7; lectura visual aproximada.\\
Refinamiento multimodal paralelo & PDF de junio, p. 8; capacidad declarada.\\
Fusión y exclusiones web & \code{app.py}, \code{judge.py}.\\
Contrato y fallback & \code{router.py}, \code{api_prompt.txt}.\\
Grafo inicial y semántico & Originales de ambos scripts en \code{experiments/}.\\
Cambios y tareas adicionales & Secciones~\ref{sec:evolucion} y~\ref{sec:extras}.\\
258 canciones alineadas & Inspección de \code{output/}.\\
\bottomrule
\end{longtable}
\normalsize
\begin{thebibliography}{9}
\bibitem{mayo} Rubén Carvajal, Leandro Chamorro y Alexandra Mendoza.
\emph{SonicSearch: Sistema RAG avanzado para recomendación musical semántica}.
OII465-1, 11 de mayo de 2026. Archivo suministrado: \code{SonicSearch.pdf}, 10 páginas.
\bibitem{junio} Rubén Carvajal, Leandro Chamorro y Alexandra Mendoza.
\emph{SonicSearch AI: Tarea 3 Módulo 4}. OII465-1, 14 de junio de 2026.
Archivo suministrado: \code{Copia de SonicSearch_IA Tarea 3 Modulo 4.pdf}, 9 páginas.
\bibitem{codigo} Proyecto SPELLFIND. Código y artefactos locales de la carpeta \code{SPELLFIND}.
Inspección del 18 de octubre de 2026; sin revisión Git identificada para fijar el snapshot.
\bibitem{enunciado} Enunciado del buscador semántico híbrido de música proporcionado por el usuario en esta conversación.
\end{thebibliography}
\end{document}
